\documentclass[../../main.tex]{subfiles}% 注意这里的文件路径不能用 ./main.tex ,否则用latexmk编译子文件会报错
\graphicspath{{\subfix{./image/}}{\subfix{../../image/}}} % 指定图片引用路径,后续可以直接使用图片文件名
% 如果图片存放在上级目录,则使用 ../../image/ 作为备用路径

\begin{document}

\section{可测函数}\label{section:ch02-可测函数}

本节建立可测函数的基本理论。可测性使可测集的原像仍为可测集，因而使简单函数的积分、可测函数的积分与像测度有定义。\cref{subsection:ch02-定义与一般性质}给出可测性的定义、生成族判据、逐点运算的封闭性、几乎处处相等与几乎处处定义的函数，以及简单函数逼近；\cref{subsection:ch02-收敛结果}比较依测度收敛、几乎一致收敛与几乎处处收敛，并证明 Egorov 定理；\cref{subsection:ch02-拟连续}讨论拟连续函数、Lusin 定理及其容量版本。

\subsection{定义与一般性质}\label{subsection:ch02-定义与一般性质}

\begin{definition}\label{definition:ch02-可测函数}\label{definition:ch02-可测性}
设 $(X,\mathscr A)$ 与 $(X',\mathscr A')$ 为可测空间，其中 $\mathscr A$ 与 $\mathscr A'$ 分别为 $X$ 与 $X'$ 上的 $\sigma$-代数。称映射 $f:X\to X'$ 为 $(\mathscr A,\mathscr A')$-可测（或简称可测），若对任意 $E'\in\mathscr A'$，其原像
\begin{align*}
f^{-1}(E')\triangleq\{x\in X\mid f(x)\in E'\}
\end{align*}
属于 $\mathscr A$。该定义只涉及两个 $\sigma$-代数，不涉及测度或拓扑。
\end{definition}

% SUPPLEMENT: parent=definition:ch02-可测函数; kind=motivation; 说明可测性的引入动机、两个 $\sigma$-代数各自的作用、原像形式的必要性以及非可测函数的边界例
% SUPPLEMENT: parent=definition:ch02-可测性; kind=motivation; 别名标签，与 definition:ch02-可测函数 共用下一条注记
\begin{remark}\phantomsection\label{remark:ch02-Borel可测}
\begin{enumerate}[(i)]
\item 引入动机。测度只定义在 $\sigma$-代数上。映射 $f:X\to X'$ 要与 $X'$ 上的测度相容，就必须把 $X'$ 中的可测集拉回为 $X$ 中的可测集；条件 $f^{-1}(\mathscr A')\subseteq\mathscr A$ 正是这一要求的精确表述。它保证对每个 $E'\in\mathscr A'$，像测度 $\mu_f(E')\triangleq\mu(f^{-1}(E'))$ 有定义，从而可以把 $X'$ 上的积分化为 $X$ 上的积分。
\item 两个 $\sigma$-代数的作用。可测性总是相对于一对 $\sigma$-代数而言：$\mathscr A$ 决定论域中哪些集合可测，$\mathscr A'$ 决定陪域中哪些集合可测。取 $\mathscr A'=\mathscr P(X')$ 或 $\mathscr A=\mathscr P(X)$ 时任何映射均可测；取 $\mathscr A=\{\varnothing,X\}$ 时，若 $\mathscr A'$ 分离 $X'$ 中的点（例如 $X'$ 为度量空间、$\mathscr A'=\mathscr B(X')$），则 $(\mathscr A,\mathscr A')$-可测函数必为常值。同一个映射相对于不同的 $\sigma$-代数可测性可以不同，故记号中必须写明这一对 $\sigma$-代数。
\item 原像形式的必要性。原像运算保持并、交、补，故可测性只需在一族生成元上验证（\cref{lemma:ch02-生成集可测}）；像运算不保持这些运算，因此定义中的原像不能换成像。
\item Borel 约定。若 $X'$ 为拓扑空间，则取 $\mathscr A'$ 为其 Borel $\sigma$-代数 $\mathscr B(X')$。此时称 $f:X\to X'$ 为 $\mathscr A$-可测（或简称可测），若它是 $(\mathscr A,\mathscr B(X'))$-可测；特别地，称 $f:X\to\mathbb{R}^N$ 为 $\mathscr A$-可测，若它是 $(\mathscr A,\mathscr B(\mathbb{R}^N))$-可测。当 $X$ 也是拓扑空间且 $\mathscr A=\mathscr B(X)$ 时，得到 Borel 函数（\cref{definition:ch02-Borel函数}）。
\item 边界情形。(a) 不可测函数存在：取 $X=X'=\mathbb{R}$、$\mathscr A=\mathscr A'=\mathscr B(\mathbb{R})$ 与 Vitali 集 $V\subseteq\mathbb{R}$，则 $V\notin\mathscr B(\mathbb{R})$，故 $\chi_V$ 不是 Borel 函数，因为 $\chi_V^{-1}(\{1\})=V$。(b) 常值函数总可测：若 $f\equiv c$，则对每个 $E'\in\mathscr A'$，$f^{-1}(E')$ 等于 $X$ 或 $\varnothing$。(c) 可测性不对像封闭：即使 $f$ 可测，集合 $f(E)$ 也未必可测，故可测性只能由原像刻画。
\end{enumerate}
\end{remark}

\begin{proposition}\label{proposition:ch02-复合可测}
设 $(X,\mathscr A)$、$(X',\mathscr A')$、$(X'',\mathscr A'')$ 为可测空间。设 $f:X\to X'$ 为 $(\mathscr A,\mathscr A')$-可测，$g:X'\to X''$ 为 $(\mathscr A',\mathscr A'')$-可测。则复合 $g\circ f:X\to X''$ 为 $(\mathscr A,\mathscr A'')$-可测。
\end{proposition}

\begin{proof}
设 $E''\in\mathscr A''$。由 $g$ 为 $(\mathscr A',\mathscr A'')$-可测得 $g^{-1}(E'')\in\mathscr A'$；再由 $f$ 为 $(\mathscr A,\mathscr A')$-可测得
\begin{align*}
(g\circ f)^{-1}(E'')=f^{-1}\bigl(g^{-1}(E'')\bigr)\in\mathscr A.
\end{align*}
由于 $E''\in\mathscr A''$ 任意，$g\circ f$ 为 $(\mathscr A,\mathscr A'')$-可测。
\end{proof}

% SUPPLEMENT: parent=proposition:ch02-复合可测; kind=motivation; 说明复合封闭性的用途、中间 $\sigma$-代数必须匹配的作用，以及不能由复合可测反推因子可测的反例
\begin{remark}\phantomsection\label{remark:ch02-复合可测}
\begin{enumerate}[(i)]
\item 引入动机。可测性对复合封闭，因而可测函数在复合下构成传递结构：连续函数与可测函数复合仍可测，可测函数经连续映射作用后仍可测。后一事实是\cref{corollary:ch02-可测运算}证明中处理 $f^2$、$|f|$ 等运算的依据。
\item 中间 $\sigma$-代数必须匹配。证明只用到 $g^{-1}(E'')\in\mathscr A'$ 与 $f^{-1}(\mathscr A')\subseteq\mathscr A$。若 $g$ 只关于比 $\mathscr A'$ 更小的 $\sigma$-代数可测，则 $g^{-1}(E'')$ 可能不属于 $\mathscr A'$，论证失效。反例：取 $X=X'=\mathbb{R}$、$\mathscr A=\mathscr A''=\mathscr B(\mathbb{R})$、$\mathscr A'=\{\varnothing,\mathbb{R}\}$，令 $f=\mathrm{id}_X$、$g=\chi_V$，其中 $V$ 为 Vitali 集。则 $f$ 为 $(\mathscr A,\mathscr A')$-可测，因为 $\varnothing$ 与 $\mathbb{R}$ 的原像仍为 $\varnothing$ 与 $\mathbb{R}$；但 $g$ 不是 $(\mathscr A',\mathscr A'')$-可测，因为 $g^{-1}(\{1\})=V\notin\mathscr A'$。此时 $g\circ f=\chi_V$ 不是 $(\mathscr A,\mathscr A'')$-可测。
\item 不能反推因子可测。由 $g\circ f$ 可测推不出 $f$ 或 $g$ 可测：取 $g\equiv0$，则 $g\circ f\equiv0$ 可测，而 $f$ 可任意（甚至不可测）。又若 $h$ 可测且 $g$ 在 $f(X)$ 上等于 $h$、在 $f(X)$ 外任意取值，则 $g\circ f=h\circ f$ 可测，而 $g$ 本身未必可测。
\item 适用边界。结论对任意可测空间成立，不涉及拓扑或测度；若把可测性换成连续性或 Borel 可测性，相应的复合封闭性仍成立，但反推的失败情形同样存在。
\end{enumerate}
\end{remark}

\begin{lemma}\label{lemma:ch02-生成集可测}
\label{lemma:ch02-基判定}
设 $(X,\mathscr A)$ 与 $(X',\mathscr A')$ 为可测空间。设 $\mathscr C'\subseteq\mathscr P(X')$ 满足极小 $\sigma$-代数 $\sigma_0(\mathscr C')=\mathscr A'$。则 $f:X\to X'$ 为 $(\mathscr A,\mathscr A')$-可测当且仅当
\begin{align}
f^{-1}(E')\in\mathscr A\quad\text{对任意 }E'\in\mathscr C'. \label{align:ch02-2.1}
\end{align}
\end{lemma}

\begin{proof}
先设 $f$ 为 $(\mathscr A,\mathscr A')$-可测，则对任意 $E'\in\mathscr A'$ 有 $f^{-1}(E')\in\mathscr A$。因 $\mathscr C'\subseteq\sigma_0(\mathscr C')=\mathscr A'$，\eqref{align:ch02-2.1}成立。

反之，设\eqref{align:ch02-2.1}成立，令
\begin{align*}
\Sigma\triangleq \{E'\subseteq X'\mid f^{-1}(E')\in\mathscr A\}.
\end{align*}
以下验证 $\Sigma$ 为 $\sigma$-代数。首先，由
\begin{align}
f^{-1}(X')=X,\quad f^{-1}(X'\setminus E')=X\setminus f^{-1}(E'),\quad
f^{-1}\!\left(\bigcup_{k=1}^\infty E'_k\right)=\bigcup_{k=1}^\infty f^{-1}(E'_k) \label{align:ch02-2.2}
\end{align}
得 $X'\in\Sigma$；若 $E'\in\Sigma$，则 $X'\setminus E'\in\Sigma$；若 $\{E'_k\}\subseteq\Sigma$，则 $\bigcup_{k=1}^\infty E'_k\in\Sigma$。故 $\Sigma$ 是 $X'$ 上的 $\sigma$-代数。另一方面，\eqref{align:ch02-2.1}等价于 $\mathscr C'\subseteq\Sigma$，于是由 $\sigma_0(\mathscr C')$ 的极小性得
\begin{align*}
\mathscr A'=\sigma_0(\mathscr C')\subseteq\Sigma,
\end{align*}
即对任意 $E'\in\mathscr A'$ 有 $f^{-1}(E')\in\mathscr A$。故 $f$ 为 $(\mathscr A,\mathscr A')$-可测。
\end{proof}

% SUPPLEMENT: parent=lemma:ch02-生成集可测; kind=motivation; 说明生成族判据的用途、条件 $\sigma_0(\mathscr C')=\mathscr A'$ 不可弱化的反例，以及生成族可取任意族的边界
% SUPPLEMENT: parent=lemma:ch02-基判定; kind=motivation; 别名标签，与 lemma:ch02-生成集可测 共用下一条注记
\begin{remark}\phantomsection\label{remark:ch02-生成族}
\begin{enumerate}[(i)]
\item 引入动机。本引理把对整个 $\sigma$-代数 $\mathscr A'$ 的验证化归到对生成族 $\mathscr C'$ 的验证。\cref{proposition:ch02-可测判据}的半直线判据、\cref{proposition:ch02-乘积可测分量}的分量判据都是它的直接应用，因而它是本节后续论证的技术基础。
\item 条件 $\sigma_0(\mathscr C')=\mathscr A'$ 的作用。该条件不可弱化为 $\mathscr C'\subseteq\mathscr A'$。反例：取 $X=X'=\mathbb{R}$、$\mathscr A=\mathscr A'=\mathscr B(\mathbb{R})$、$\mathscr C'=\{\varnothing,\mathbb{R}\}$，则 $\sigma_0(\mathscr C')=\{\varnothing,\mathbb{R}\}\neq\mathscr A'$；函数 $\chi_V$（$V$ 为 Vitali 集）满足对任意 $E'\in\mathscr C'$ 有 $\chi_V^{-1}(E')\in\mathscr A$，但它不是 Borel 函数。原因是 $\Sigma$ 只保证包含 $\mathscr C'$ 生成的 $\sigma$-代数，而不保证包含更大的 $\mathscr A'$。
\item 适用边界。$\mathscr C'$ 可为任意族：不必是 $\sigma$-代数，不必可数，也不必含 $X'$。条件只通过 $\sigma_0(\mathscr C')$ 起作用，故把 $\mathscr C'$ 换成任何生成同一 $\sigma$-代数的族，结论不变。当 $\mathscr C'=\varnothing$ 时 $\sigma_0(\mathscr C')=\{\varnothing,X'\}$，条件空洞成立，此时任何 $f:X\to X'$ 均 $(\mathscr A,\{\varnothing,X'\})$-可测，与定义一致。
\end{enumerate}
\end{remark}

\begin{proposition}\label{proposition:ch02-限制可测}\label{proposition:ch02-可测拼接}
设 $(X,\mathscr A)$ 与 $(X',\mathscr A')$ 为可测空间，$X=\bigcup_{k=1}^\infty E_k$，$\{E_k\}\subseteq\mathscr A$。则 $f:X\to X'$ 为 $(\mathscr A,\mathscr A')$-可测当且仅当每个限制 $f|_{E_k}:E_k\to X'$ 为 $(\mathscr A\cap E_k,\mathscr A')$-可测，其中 $\mathscr A\cap E_k$ 为 $\mathscr A$ 在 $E_k$ 上的迹 $\sigma$-代数。
\end{proposition}

\begin{proof}
设对每个 $k\in\mathbb{N}$，$f|_{E_k}$ 为 $(\mathscr A\cap E_k,\mathscr A')$-可测。对任意 $F'\in\mathscr A'$，有 $(f|_{E_k})^{-1}(F')=E_k\cap f^{-1}(F')\in\mathscr A\cap E_k\subseteq\mathscr A$，这里用到 $E_k\in\mathscr A$。于是
\begin{align*}
f^{-1}(F')=\left(\bigcup_{k=1}^\infty E_k\right)\cap f^{-1}(F')
=\bigcup_{k=1}^\infty\bigl(E_k\cap f^{-1}(F')\bigr)
=\bigcup_{k=1}^\infty (f|_{E_k})^{-1}(F')\in\mathscr A,
\end{align*}
从而 $f$ 为 $(\mathscr A,\mathscr A')$-可测。

反之，设 $f$ 为 $(\mathscr A,\mathscr A')$-可测。注意 $f|_{E_k}=f\circ j_{E_k}$，其中 $j_{E_k}:E_k\to X$ 为嵌入。对任意 $E\in\mathscr A$，
\begin{align*}
j_{E_k}^{-1}(E)\triangleq \{x\in E_k\mid j_{E_k}(x)\in E\}=E\cap E_k\in\mathscr A\cap E_k,
\end{align*}
故 $j_{E_k}$ 为 $(\mathscr A\cap E_k,\mathscr A)$-可测。由\cref{proposition:ch02-复合可测}得 $f|_{E_k}$ 为 $(\mathscr A\cap E_k,\mathscr A')$-可测（$k\in\mathbb{N}$）。
\end{proof}

% SUPPLEMENT: parent=proposition:ch02-限制可测; kind=motivation; 说明逐块验证可测性的用途、$E_k\in\mathscr A$ 不可省的反例，以及可数覆盖中可数性的作用
% SUPPLEMENT: parent=proposition:ch02-可测拼接; kind=motivation; 别名标签，与 proposition:ch02-限制可测 共用下一条注记
\begin{remark}\phantomsection\label{remark:ch02-限制可测}
\begin{enumerate}[(i)]
\item 引入动机。本命题把整体可测性化归为可数多个可测块上的可测性，是拼接论证的标准工具：例如在 $\mathbb{R}^N$ 上按球层 $B(0,k)\setminus B(0,k-1)$ 分块，或把区域分解为可数多个可测部分时，只需逐块验证。
\item $E_k\in\mathscr A$ 的作用。该假设在正反两个方向都不可省。为看"当"方向，取 $X=\{0,1\}$、$\mathscr A=\{\varnothing,X\}$、$E=\{0\}\notin\mathscr A$、$f=\chi_{\{0\}}$；则 $f|_E$ 为常值 $1$，故 $(\mathscr A\cap E,\mathscr A')$-可测，但 $f^{-1}(\{1\})=\{0\}\notin\mathscr A$，$f$ 不可测。为看"仅当"方向，若 $E_k\notin\mathscr A$，则迹 $\sigma$-代数 $\mathscr A\cap E_k$ 仍可定义，但嵌入 $j_{E_k}$ 未必可测，上述复合论证失效。
\item 可数性的作用。证明中只用到可数并，因为 $\sigma$-代数对可数并封闭。对不可数覆盖，即使每个限制可测，把 $f^{-1}(F')$ 写成不可数并也得不到可测性。有限覆盖是本命题的特例：把 $E_k$ 换成 $E_k\cup\varnothing$ 即可补足为可数族。
\item 适用边界。若 $E_k$ 两两不交且覆盖 $X$，结论给出 $f$ 的可测性与各分块上可测性的等价；若只覆盖 $X$ 的一个子集，则只能得到 $f$ 在该子集上的可测性。
\end{enumerate}
\end{remark}

对任意 $f:X\to\mathbb{R}$ 与 $\alpha\in\mathbb{R}$，记 $\{f>\alpha\}\triangleq\{x\in X\mid f(x)>\alpha\}=f^{-1}((\alpha,\infty))$；集合 $\{f\geqslant\alpha\}$、$\{f<\alpha\}$、$\{f\leqslant\alpha\}$、$\{f=\alpha\}$ 按同样方式定义。对实值函数有如下可测性判据。

\begin{proposition}\label{proposition:ch02-可测判据}\label{proposition:ch02-Borel判定}
设 $(X,\mathscr A)$ 为可测空间，$f:X\to\mathbb{R}$ 为函数。则下列命题等价：
\begin{enumerate}[(i)]
\item $f$ 为 $\mathscr A$-可测；
\item $\{f>\alpha\}\in\mathscr A$ 对任意 $\alpha\in\mathbb{R}$；
\item $\{f\geqslant\alpha\}\in\mathscr A$ 对任意 $\alpha\in\mathbb{R}$；
\item $\{f<\alpha\}\in\mathscr A$ 对任意 $\alpha\in\mathbb{R}$；
\item $\{f\leqslant\alpha\}\in\mathscr A$ 对任意 $\alpha\in\mathbb{R}$。
\end{enumerate}
\end{proposition}

% SUPPLEMENT: parent=proposition:ch02-可测判据; kind=proof; 原书由生成族引理直接断言，此处补出 (ii) 与 (iii)、(iv) 与 (v) 之间用可数交转换的完整等价链
% SUPPLEMENT: parent=proposition:ch02-Borel判定; kind=proof; 别名标签，与 proposition:ch02-可测判据 共用同一补证
\begin{proof}
记 $\mathscr H\triangleq\{(\alpha,\infty)\mid\alpha\in\mathbb{R}\}$。由\hyperref[remark:ch01-实直线Borel]{实直线 Borel $\sigma$-代数的生成族}得 $\sigma_0(\mathscr H)=\mathscr B(\mathbb{R})$。

(i)$\Rightarrow$(ii)：由 $\mathscr H\subseteq\mathscr B(\mathbb{R})$ 与 $\mathscr A$-可测性的定义，对任意 $\alpha\in\mathbb{R}$ 有 $\{f>\alpha\}=f^{-1}((\alpha,\infty))\in\mathscr A$。

(ii)$\Rightarrow$(i)：由\cref{lemma:ch02-生成集可测}，只需对 $\mathscr C'=\mathscr H$ 验证\eqref{align:ch02-2.1}；而 (ii) 恰断言 $f^{-1}(H)\in\mathscr A$ 对每个 $H\in\mathscr H$ 成立，故 $f$ 为 $\mathscr A$-可测。

(ii)$\Rightarrow$(iii)：对 $\alpha\in\mathbb{R}$ 有
\begin{align*}
\{f\geqslant\alpha\}=\bigcap_{n=1}^\infty\{f>\alpha-1/n\},
\end{align*}
因为 $f(x)\geqslant\alpha$ 当且仅当对每个 $n\in\mathbb{N}$ 有 $f(x)>\alpha-1/n$。故 (iii) 成立。

(iii)$\Rightarrow$(iv)：对 $\alpha\in\mathbb{R}$ 有 $\{f<\alpha\}=X\setminus\{f\geqslant\alpha\}\in\mathscr A$。

(iv)$\Rightarrow$(v)：对 $\alpha\in\mathbb{R}$ 有
\begin{align*}
\{f\leqslant\alpha\}=\bigcap_{n=1}^\infty\{f<\alpha+1/n\},
\end{align*}
因为 $f(x)\leqslant\alpha$ 当且仅当对每个 $n\in\mathbb{N}$ 有 $f(x)<\alpha+1/n$。故 (v) 成立。

(v)$\Rightarrow$(ii)：对 $\alpha\in\mathbb{R}$ 有 $\{f>\alpha\}=X\setminus\{f\leqslant\alpha\}\in\mathscr A$。

于是 (i)$\Rightarrow$(ii)$\Rightarrow$(iii)$\Rightarrow$(iv)$\Rightarrow$(v)$\Rightarrow$(ii) 且 (ii)$\Rightarrow$(i)，五个命题等价。
\end{proof}

% SUPPLEMENT: parent=proposition:ch02-可测判据; kind=motivation; 说明半直线判据的用途、可数交在等价链中的作用、只检验部分 $\alpha$ 不足的反例，以及扩充实值情形的边界
% SUPPLEMENT: parent=proposition:ch02-Borel判定; kind=motivation; 别名标签，与 proposition:ch02-可测判据 共用下一条注记
\begin{remark}\phantomsection\label{remark:ch02-可测判据}
\begin{enumerate}[(i)]
\item 引入动机。本命题把"对每个 Borel 集验证原像可测"化归为"对每个半直线验证原像可测"，使可测性成为可以用水平集检验的性质。\cref{corollary:ch02-可测运算}以及简单函数逼近定理都反复使用这一判据。
\item 条件的作用。$\alpha$ 必须取遍整个 $\mathbb{R}$。(ii) 与 (iii)、(iv) 与 (v) 之间的转换用到可数交，其可行性依赖于 $\alpha-1/n\downarrow\alpha$ 与 $\alpha+1/n\uparrow\alpha$；这正是 $\sigma$-代数对可数运算封闭之处。只对有限个或可数个 $\alpha$ 检验不足以推出可测性：取 $\mathscr A=\{\varnothing,X\}$、$f=\chi_V$（$V$ 为 Vitali 集），则 $\{f>-1\}=X$、$\{f>2\}=\varnothing$ 均属于 $\mathscr A$，但 $f$ 不可测。
\item 适用边界。对扩充实值函数 $f:X\to[-\infty,\infty]$，半直线族 $\{(\alpha,\infty]\mid\alpha\in\mathbb{R}\}$ 生成 $\mathscr B([-\infty,\infty])$，故同一判据在扩充实值 Borel $\sigma$-代数下成立。判据中的"开半直线"也可换成闭半直线、开区间或闭区间等任一生成 $\mathscr B(\mathbb{R})$ 的族。
\end{enumerate}
\end{remark}

\begin{corollary}\label{corollary:ch02-可测运算}\label{corollary:ch02-逐点极限}
设 $(X,\mathscr A)$ 为可测空间。
\begin{enumerate}[(i)]
\item 若 $f,g:X\to\mathbb{R}$ 为 $\mathscr A$-可测，则函数 $f+g$、$fg$、$|f|$、$\min\{f,g\}$、$\max\{f,g\}$ 以及 $f/g$（后者要求 $g$ 在 $X$ 上处处非零）均为 $\mathscr A$-可测。
\item 若对每个 $n\in\mathbb{N}$，函数 $f_n:X\to[-\infty,\infty]$ 为 $\mathscr A$-可测，则函数 $\sup_{n\in\mathbb{N}}f_n$、$\inf_{n\in\mathbb{N}}f_n$、$\limsup_{n\to\infty}f_n$、$\liminf_{n\to\infty}f_n$ 均为 $\mathscr A$-可测。特别地，若逐点极限 $\lim_{n\to\infty}f_n$ 存在，则它为 $\mathscr A$-可测。
\end{enumerate}
\end{corollary}

% SUPPLEMENT: parent=corollary:ch02-可测运算; kind=proof; 原书以"由此立得"略去，补出连续函数复合、加法用 $\mathbb{Q}$ 的稠密性、乘除与格运算的化归，以及扩充实值上确界与上下极限的完整论证
% SUPPLEMENT: parent=corollary:ch02-逐点极限; kind=proof; 别名标签，与 corollary:ch02-可测运算 共用同一补证
\begin{proof}
先证一个事实：若 $\varphi:\mathbb{R}\to\mathbb{R}$ 连续、$f:X\to\mathbb{R}$ 为 $\mathscr A$-可测，则 $\varphi\circ f$ 为 $\mathscr A$-可测。事实上，对任意 $B\in\mathscr B(\mathbb{R})$，由 $\varphi$ 连续得 $\varphi^{-1}(B)\in\mathscr B(\mathbb{R})$，故
\begin{align*}
(\varphi\circ f)^{-1}(B)=f^{-1}\bigl(\varphi^{-1}(B)\bigr)\in\mathscr A.
\end{align*}

(i) 加法。对 $\alpha\in\mathbb{R}$ 有
\begin{align*}
\{f+g>\alpha\}=\bigcup_{q\in\mathbb{Q}}\bigl(\{f>q\}\cap\{g>\alpha-q\}\bigr).
\end{align*}
若 $f(x)+g(x)>\alpha$，则区间 $(\alpha-g(x),f(x))$ 非空，取有理数 $q\in(\alpha-g(x),f(x))$，得 $f(x)>q$ 且 $g(x)>\alpha-q$；反之，若存在 $q\in\mathbb{Q}$ 使 $f(x)>q$、$g(x)>\alpha-q$，则 $f(x)+g(x)>\alpha$。由于 $\mathbb{Q}$ 可数，上式右端为可数并，故属于 $\mathscr A$；由\cref{proposition:ch02-可测判据}得 $f+g$ 可测。

平方与绝对值。取 $\varphi(t)=t^2$、$\varphi(t)=|t|$，由上述事实得 $f^2$ 与 $|f|$ 可测。

乘积。由恒等式
\begin{align*}
fg=\frac{(f+g)^2-f^2-g^2}{2}
\end{align*}
及已证的加法与数乘（数乘为与连续函数 $t\mapsto ct$ 的复合）得 $fg$ 可测。

格运算。由恒等式
\begin{align*}
\min\{f,g\}=\frac{f+g-|f-g|}{2},\qquad
\max\{f,g\}=\frac{f+g+|f-g|}{2}
\end{align*}
得 $\min\{f,g\}$ 与 $\max\{f,g\}$ 可测。

倒数。设 $g$ 在 $X$ 上处处非零。对 $\alpha\in\mathbb{R}$ 分三种情形：
\begin{align*}
\{1/g>\alpha\}=
\begin{cases}
\{0<g<1/\alpha\}, & \alpha>0,\\
\{g>0\}, & \alpha=0,\\
\{g>0\}\cup\{g<1/\alpha\}, & \alpha<0.
\end{cases}
\end{align*}
三种情形右端均属于 $\mathscr A$，故由\cref{proposition:ch02-可测判据}得 $1/g$ 可测。再由乘积的可测性得 $f/g=f\cdot(1/g)$ 可测。

(ii) 先注意：对任意 $h:X\to[-\infty,\infty]$，$h$ 为 $\mathscr A$-可测当且仅当 $\{h>\alpha\}\in\mathscr A$ 对任意 $\alpha\in\mathbb{R}$ 成立。这是因为 $\mathscr B([-\infty,\infty])=\sigma_0(\{(\alpha,\infty]\mid\alpha\in\mathbb{R}\})$，再用\cref{lemma:ch02-生成集可测}即可。

上确界与下确界。对 $\alpha\in\mathbb{R}$ 有
\begin{align*}
\left\{\sup_{n\in\mathbb{N}}f_n>\alpha\right\}=\bigcup_{n=1}^\infty\{f_n>\alpha\}\in\mathscr A,
\qquad
\left\{\inf_{n\in\mathbb{N}}f_n<\alpha\right\}=\bigcup_{n=1}^\infty\{f_n<\alpha\}\in\mathscr A,
\end{align*}
故 $\sup_{n\in\mathbb{N}}f_n$ 与 $\inf_{n\in\mathbb{N}}f_n$ 可测。

上极限与下极限。对每个 $m\in\mathbb{N}$，函数 $g_m\triangleq\sup_{n\geqslant m}f_n$ 与 $h_m\triangleq\inf_{n\geqslant m}f_n$ 可测（对上确界与下确界的结论应用于序列 $n\mapsto f_{n+m-1}$）。于是
\begin{align*}
\limsup_{n\to\infty}f_n=\inf_{m\in\mathbb{N}}g_m,\qquad
\liminf_{n\to\infty}f_n=\sup_{m\in\mathbb{N}}h_m
\end{align*}
可测。若逐点极限存在，则它等于 $\limsup_{n\to\infty}f_n$，故可测。
\end{proof}

% SUPPLEMENT: parent=corollary:ch02-可测运算; kind=motivation; 说明运算封闭性的用途、$g\neq0$ 与实值性的作用、不可数上确界失效的反例以及 $\pm\infty$ 运算的边界
% SUPPLEMENT: parent=corollary:ch02-逐点极限; kind=motivation; 别名标签，与 corollary:ch02-可测运算 共用下一条注记
\begin{remark}\phantomsection\label{remark:ch02-可测运算}
\begin{enumerate}[(i)]
\item 引入动机。可测函数类对四则运算、绝对值、格运算与可数逐点极限封闭，这是积分理论的技术前提：单调收敛定理与 Fatou 引理中的 $\sup_n f_n$、$\liminf_n f_n$ 必须可测，控制收敛定理中的 $|f_n-f|$ 也必须可测。
\item 条件的作用。$g$ 在 $X$ 上处处非零是 $f/g$ 作为 $X$ 上函数的定义条件；若只假设 $g\neq0$ $\mu$-a.e.，则 $f/g$ 只在 $\mu$-a.e. 意义下有定义，其可测性须按\cref{definition:ch02-a.e.可测函数}理解。(i) 只对实值函数陈述，因为 $[-\infty,\infty]$ 值函数的和与积可能出现 $\infty-\infty$、$0\cdot\infty$ 这类无定义情形；(ii) 中的上确界、下确界与上下极限则对扩充实值函数总有定义。
\item 不可数上确界失效。设 $A\subseteq\mathbb{R}$ 不是 Borel 集，对每个 $x\in A$ 令 $f_x\triangleq\chi_{\{x\}}$；每个 $f_x$ 为 Borel 可测，但 $\sup_{x\in A}f_x=\chi_A$ 不可测。故 (ii) 中的"可数"不能换成"不可数"。
\item 适用边界。逐点极限存在时才断言极限可测；若极限不存在，$\limsup$ 与 $\liminf$ 仍可测，且二者之差刻画了振荡。上述证明只用到 $\mathbb{Q}$ 的可数性与 $\sigma$-代数对可数并、交、补的封闭性，不依赖测度。
\end{enumerate}
\end{remark}

\begin{proposition}\label{proposition:ch02-a.e.相等可测}
设 $(X,\mathscr A,\mu)$ 为完备测度空间，$(X',\mathscr A')$ 为可测空间。设 $f,g:X\to X'$ 在 $X$ 上 $\mu$-a.e. 相等。则 $f$ 为 $(\mathscr A,\mathscr A')$-可测当且仅当 $g$ 为 $(\mathscr A,\mathscr A')$-可测。
\end{proposition}

\begin{proof}
由假设存在零测集 $N\in\mathscr A$（即 $\mu(N)=0$），使对任意 $x\in N^c$ 有 $g(x)=f(x)$。对任意 $E'\in\mathscr A'$，有集合等式
\begin{align*}
g^{-1}(E')=\bigl(g^{-1}(E')\cap N\bigr)\cup\bigl(f^{-1}(E')\cap N^c\bigr).
\end{align*}
事实上，右端含于左端：第一项含于 $g^{-1}(E')$；第二项中 $x\in N^c$ 且 $f(x)\in E'$，故 $g(x)=f(x)\in E'$。反之，若 $x\in g^{-1}(E')$，则 $x\in N$ 或 $x\in N^c$；前者属于第一项，后者满足 $f(x)=g(x)\in E'$，属于第二项。

若 $f$ 为 $(\mathscr A,\mathscr A')$-可测，则 $f^{-1}(E')\cap N^c\in\mathscr A$。又 $g^{-1}(E')\cap N\subseteq N$ 且 $\mu(N)=0$，由 $\mu$ 的完备性得 $g^{-1}(E')\cap N\in\mathscr A$。于是 $g^{-1}(E')\in\mathscr A$，即 $g$ 为 $(\mathscr A,\mathscr A')$-可测。交换 $f$ 与 $g$ 的角色即得另一方向。
\end{proof}

% SUPPLEMENT: parent=proposition:ch02-a.e.相等可测; kind=motivation; 说明完备性的作用、去掉完备性后的反例、几乎处处相等作为等价关系的意义，以及向几乎处处定义函数推广的边界
\begin{remark}\phantomsection\label{remark:ch02-a.e.相等}
\begin{enumerate}[(i)]
\item 引入动机。完备化把零测集的子集也纳入 $\sigma$-代数，使"几乎处处相等"成为可测函数之间的等价关系。于是可以把只在零测集上不同的函数视为同一对象，这是 $L^0$ 空间与 $L^p$ 空间构造的起点。
\item 完备性不可省。反例：取 $X=\{0,1,2\}$、$\mathscr A=\{\varnothing,\{0\},\{1,2\},X\}$，令 $\mu(\{0\})=1$、$\mu(\{1,2\})=0$。此时 $\{1,2\}$ 为零测集，但其子集 $\{1\}\notin\mathscr A$，故 $\mu$ 不完备。令 $f\equiv0$、$g=\chi_{\{1\}}$。则 $f$ 为 $\mathscr A$-可测，$f=g$ 在 $X\setminus\{1\}$ 上成立，而 $\{f\neq g\}=\{1\}\subseteq\{1,2\}$ 且 $\mu(\{1,2\})=0$，但 $g^{-1}(\{1\})=\{1\}\notin\mathscr A$，故 $g$ 不可测。证明中 $g^{-1}(E')\cap N\in\mathscr A$ 这一步正需要完备性。
\item 与几乎处处定义的函数的关系。若 $f,g$ 只在某个零测集之外有定义并在那里 $\mu$-a.e. 相等，则同样的结论在\cref{definition:ch02-a.e.可测函数}的意义下成立：把测度空间限制到公共定义域上，完备性被迹 $\sigma$-代数继承，再应用本命题即可。
\item 适用边界。命题对任意陪域可测空间成立，不要求 $X'$ 为拓扑空间；若把"$\mu$-a.e. 相等"换成"$C$-q.e. 相等"，则需把完备性换成容量的相应封闭性，结论一般不再成立。
\end{enumerate}
\end{remark}

若 $\mu$ 不完备，则\cref{proposition:ch02-a.e.相等可测}的结论不必成立（反例见\hyperref[remark:ch02-a.e.相等]{上述注记}）。为避免这一困难，将可测性定义推广到几乎处处定义的函数是方便的。

\begin{definition}\label{definition:ch02-a.e.可测函数}\label{definition:ch02-几乎处处可测}
设 $(X,\mathscr A,\mu)$ 为测度空间，$(X',\mathscr A')$ 为可测空间。设 $N\subseteq X$ 为零测集，即 $N\in\mathscr A$ 且 $\mu(N)=0$，并设 $f:N^c\to X'$ 为定义在 $N^c\triangleq X\setminus N$ 上的映射。称 $f$ 为 $(\mathscr A,\mathscr A')$-可测，若它作为 $(\mathscr A\cap N^c,\mathscr A')$-可测映射，即对任意 $E'\in\mathscr A'$ 有 $f^{-1}(E')\in\mathscr A\cap N^c$。
\end{definition}

% SUPPLEMENT: parent=definition:ch02-a.e.可测函数; kind=motivation; 说明引入几乎处处定义函数的动机、$N\in\mathscr A$ 与 $\mu(N)=0$ 各自的作用、延拓到全空间的构造与边界情形
% SUPPLEMENT: parent=definition:ch02-几乎处处可测; kind=motivation; 别名标签，与 definition:ch02-a.e.可测函数 共用下一条注记
\begin{remark}\phantomsection\label{remark:ch02-a.e.极限可测}\label{remark:ch02-几乎处处}
\begin{enumerate}[(i)]
\item 引入动机。若 $\mu$ 不完备，两个几乎处处相等的函数可以一个可测而另一个不可测（\cref{proposition:ch02-a.e.相等可测}的反例）。把"定义在几乎整个 $X$ 上"的函数直接纳入可测性框架，可以避免先完备化测度空间。
\item 条件的作用。$N\in\mathscr A$ 使迹 $\sigma$-代数 $\mathscr A\cap N^c$ 有定义；$\mu(N)=0$ 使 $N$ 可以被忽略。若 $N\notin\mathscr A$，则 $N^c\notin\mathscr A$，上述迹 $\sigma$-代数虽仍可定义，但无法用延拓论证把 $f$ 化为 $X$ 上的可测函数。当 $N=\varnothing$ 时 $\mathscr A\cap N^c=\mathscr A$，上述定义化归为\cref{definition:ch02-可测函数}。
\item 延拓到全空间。设 $X'\neq\varnothing$，取定 $x_0'\in X'$，令 $\tilde f(x)\triangleq f(x)$（$x\in N^c$）、$\tilde f(x)\triangleq x_0'$（$x\in N$）。则对任意 $E'\in\mathscr A'$，
\begin{align*}
\tilde f^{-1}(E')=f^{-1}(E')\cup(N\ \text{若}\ x_0'\in E').
\end{align*}
右端属于 $\mathscr A$，故 $\tilde f$ 为 $(\mathscr A,\mathscr A')$-可测，且 $\tilde f=f$ 在 $N^c$ 上成立。
\item 几乎处处极限可测。设 $\{f_n\}$ 为 $\mathscr A$-可测函数列，且 $f_n\to g$ 在 $X$ 上 $\mu$-a.e.。则存在零测集 $N\in\mathscr A$ 使 $f_n(x)\to g(x)$ 对每个 $x\in N^c$ 成立。把 $g$ 视为定义在 $(N^c)$ 上的函数；在迹空间 $(N^c,\mathscr A\cap N^c)$ 上，每个 $f_n|_{N^c}$ 可测，且逐点极限为 $g$。当 $X'=\mathbb{R}$ 时，由\cref{corollary:ch02-可测运算}(ii) 得 $g$ 为 $(\mathscr A\cap N^c,\mathscr A')$-可测，即 $g$ 在\cref{definition:ch02-a.e.可测函数}的意义下可测；对取值于一般可测空间的映射，逐点极限的可测性需另行假设。当 $\mu$ 完备时，由\cref{proposition:ch02-a.e.相等可测}还可把 $g$ 延拓为 $X$ 上的 $\mathscr A$-可测函数。
\item 边界情形。\cref{definition:ch02-a.e.可测函数}意义下的可测性只依赖于 $f$ 在 $N^c$ 上的取值；在零测集上改变 $f$ 不影响可测性。反之，若两个函数只在零测集上不同，则它们在\cref{definition:ch02-a.e.可测函数}的意义下同时可测或同时不可测。
\end{enumerate}
\end{remark}

下一条定义给出可测性的拓扑版本。

\begin{definition}\label{definition:ch02-Borel函数}
设 $(X,\mathscr T)$ 与 $(X',\mathscr T')$ 为拓扑空间，$\mathscr B\triangleq\sigma_0(\mathscr T)$ 与 $\mathscr B'\triangleq\sigma_0(\mathscr T')$ 为相应的 Borel $\sigma$-代数。称 $f:X\to X'$ 为 Borel 可测（或 Borel 函数），若它为 $(\mathscr B,\mathscr B')$-可测，即对任意 $E'\in\mathscr B'$ 有 $f^{-1}(E')\in\mathscr B$。
\end{definition}

% SUPPLEMENT: parent=definition:ch02-Borel函数; kind=motivation; 说明 Borel 可测性的动机、$\sigma_0(\mathscr T)$ 的作用、连续与半连续函数的可测性、与 Lebesgue 可测性的差别
\begin{remark}\phantomsection\label{remark:ch02-连续半连续Borel}
\phantomsection\label{remark:ch02-连续Borel}\label{remark:ch02-连续逼近注}
\begin{enumerate}[(i)]
\item 引入动机。拓扑空间自带开集族，由它生成的 Borel $\sigma$-代数给出最常用的可测结构。Borel 可测性是"与拓扑相容"的可测性：连续映射把开集拉回为开集，因而自动 Borel 可测。
\item 连续函数。若 $(X,\mathscr T)$ 与 $(X',\mathscr T')$ 为拓扑空间，$f:X\to X'$ 连续，则对任意 $E'\in\mathscr T'$ 有 $f^{-1}(E')\in\mathscr T\subseteq\mathscr B$。由\cref{lemma:ch02-生成集可测}（取 $\mathscr C'=\mathscr T'$，$\sigma_0(\mathscr T')=\mathscr B'$）得 $f$ 为 $(\mathscr B,\mathscr B')$-可测。
\item 半连续函数。设 $(X,\mathscr T)$ 为拓扑空间。称 $f:X\to(-\infty,\infty]$ 为下半连续，若对任意 $x\in X$ 有 $\liminf_{y\to x}f(y)\geqslant f(x)$，等价地，对任意 $\alpha\in\mathbb{R}$，集合 $\{f>\alpha\}$ 为开集。等价性的验证如下：若 $f$ 下半连续，则对 $x\in\{f>\alpha\}$ 与 $\varepsilon\triangleq(f(x)-\alpha)/2>0$，存在 $x$ 的邻域 $U$ 使 $y\in U$ 时 $f(y)>f(x)-\varepsilon>\alpha$，故 $\{f>\alpha\}$ 为开集；反之，若所有 $\{f>\alpha\}$ 为开集，则对 $x\in X$ 与 $\varepsilon>0$，$U\triangleq\{f>f(x)-\varepsilon\}$ 是含 $x$ 的开集，且 $y\in U$ 时 $f(y)>f(x)-\varepsilon$，故 $\liminf_{y\to x}f(y)\geqslant f(x)-\varepsilon$，由 $\varepsilon$ 任意得下半连续。类似地，$f:X\to[-\infty,\infty)$ 为上半连续当且仅当对任意 $\alpha\in\mathbb{R}$，$\{f<\alpha\}$ 为开集。于是由\cref{proposition:ch02-可测判据}，下半连续与上半连续函数均为 Borel 可测。
\item 与 Lebesgue 可测性的差别。设 $X=X'=\mathbb{R}$。Borel $\sigma$-代数 $\mathscr B(\mathbb{R})$ 严格包含于 Lebesgue $\sigma$-代数 $\mathscr L$，故存在 $\mathscr L$-可测而 $\mathscr B(\mathbb{R})$-不可测的函数（例如 $\chi_V$，其中 $V\in\mathscr L\setminus\mathscr B(\mathbb{R})$）；反之，每个 Borel 函数都是 Lebesgue 可测的。本书如无特别说明，$\mathbb{R}^N$ 上的"可测"均指 Borel 可测。
\end{enumerate}
\end{remark}

\begin{definition}\label{definition:ch02-简单函数}
设 $(X,\mathscr A)$ 为可测空间。称函数 $s:X\to\mathbb{R}$ 为简单函数，若存在 $n\in\mathbb{N}$、实数 $c_1,\cdots,c_n\in\mathbb{R}$ 与 $X$ 的划分 $\{E_1,\cdots,E_n\}\subseteq\mathscr A$（即 $E_i\cap E_j=\varnothing$ 当 $i\neq j$，且 $\bigcup_{i=1}^n E_i=X$），使
\begin{align*}
s=\sum_{i=1}^n c_i\chi_{E_i}.
\end{align*}
\end{definition}

% SUPPLEMENT: parent=definition:ch02-简单函数; kind=motivation; 说明简单函数的动机、划分条件的作用、表示不唯一与规范表示、可测性以及 $\mathscr S(X)$ 与 $\mathscr S^+(X)$ 的封闭性
\begin{remark}\phantomsection\label{remark:ch02-简单函数}
\begin{enumerate}[(i)]
\item 引入动机。简单函数只取有限个值，其积分可直接由 $\sum_{i=1}^n c_i\mu(E_i)$ 定义（\cref{definition:ch02-简单函数积分}）。一般非负可测函数用非降简单函数列从下方逼近（\cref{theorem:ch02-简单逼近}），从而积分由简单函数的积分取上确界定义。
\item 划分条件的作用。要求 $\{E_i\}$ 为划分（两两不交且并为 $X$）使每个 $x\in X$ 恰属于一个 $E_i$，从而 $s(x)=c_i$ 被唯一确定，表示式不会自相矛盾。若只要求 $\{E_i\}$ 覆盖 $X$，则同一函数可能在不同 $E_i$ 上被赋予不同的值。允许 $E_i=\varnothing$ 或 $c_i$ 重复不影响 $s$；若要求 $c_i$ 互异且 $E_i=\{s=c_i\}$，则得到唯一的规范表示（至多相差指标重排与 $E_i=\varnothing$ 的项）。
\item 可测性。对任意 $B\in\mathscr B(\mathbb{R})$ 有 $s^{-1}(B)=\bigcup_{i:\,c_i\in B}E_i\in\mathscr A$，故每个简单函数均为 $\mathscr A$-可测。
\item 封闭性。记 $\mathscr S(X)$ 为实值简单函数的集合，$\mathscr S^+(X)$ 为非负简单函数 $s:X\to[0,\infty)$ 的集合。若 $s=\sum_i a_i\chi_{A_i}$、$t=\sum_j b_j\chi_{B_j}$ 为简单函数，则在公共划分 $\{A_i\cap B_j\}$ 上 $s+t=\sum_{i,j}(a_i+b_j)\chi_{A_i\cap B_j}$、$st=\sum_{i,j}a_ib_j\chi_{A_i\cap B_j}$，故 $\mathscr S(X)$ 对加法、乘法与数乘封闭。简单函数只取有限个值，故 $n=1$ 时常值函数是简单函数。
\end{enumerate}
\end{remark}

下列结果的证明是熟知的。

\begin{theorem}\label{theorem:ch02-简单逼近}\label{theorem:ch02-简单函数逼近}
设 $(X,\mathscr A)$ 为可测空间，$f:X\to\mathbb{R}$ 为 $\mathscr A$-可测函数。则存在简单函数列 $\{s_n\}$ 使 $s_n\to f$ 在 $X$ 上逐点收敛。此外：
\begin{enumerate}[(i)]
\item 若 $f\geqslant0$，则可取 $\{s_n\}$ 非降，且对任意 $n\in\mathbb{N}$ 有 $0\leqslant s_n\leqslant f$；
\item 若 $f$ 有界，则 $s_n\to f$ 在 $X$ 上一致收敛。
\end{enumerate}
\end{theorem}

% SUPPLEMENT: parent=theorem:ch02-简单逼近; kind=proof; 原书称证明熟知而未给出，此处补出二进水平集构造、非降性、逐点收敛、有界时的一致收敛估计以及一般实值函数的正负部分解
% SUPPLEMENT: parent=theorem:ch02-简单函数逼近; kind=proof; 别名标签，与 theorem:ch02-简单逼近 共用同一补证
\begin{proof}
先设 $f\geqslant0$。对 $n\in\mathbb{N}$ 与 $k=1,\cdots,n2^n$，令
\begin{align*}
E_{n,k}\triangleq\{x\in X\mid (k-1)/2^n\leqslant f(x)<k/2^n\},\qquad
F_n\triangleq\{f\geqslant n\}.
\end{align*}
由\cref{proposition:ch02-可测判据}，每个 $E_{n,k}$ 与 $F_n$ 属于 $\mathscr A$，且 $\{E_{n,1},\cdots,E_{n,n2^n},F_n\}$ 是 $X$ 的划分。令
\begin{align*}
s_n\triangleq\sum_{k=1}^{n2^n}\frac{k-1}{2^n}\chi_{E_{n,k}}+n\chi_{F_n}.
\end{align*}
则 $s_n$ 为简单函数。并且 $0\leqslant s_n\leqslant f$：在 $E_{n,k}$ 上 $s_n=(k-1)/2^n\leqslant f$，在 $F_n$ 上 $s_n=n\leqslant f$。

非降性。固定 $x\in X$ 并令 $t\triangleq f(x)$。若 $t\geqslant n+1$，则 $s_n(x)=n$ 且 $s_{n+1}(x)=n+1$。若 $n\leqslant t<n+1$，则 $s_n(x)=n$，而 $x\in E_{n+1,k}$ 对某个 $k$ 成立，且 $(k-1)/2^{n+1}\geqslant n$，故 $s_{n+1}(x)\geqslant n=s_n(x)$。若 $t<n$，写 $m\triangleq\lfloor2^nt\rfloor$，则 $s_n(x)=m/2^n$，且由 $m\leqslant2^nt<m+1$ 得 $2m\leqslant2^{n+1}t$，故 $\lfloor2^{n+1}t\rfloor\geqslant2m$，即 $s_{n+1}(x)\geqslant m/2^n=s_n(x)$。三种情形合起来给出 $s_n(x)\leqslant s_{n+1}(x)$。

逐点收敛。固定 $x\in X$。若 $t=f(x)<\infty$，取 $n>t$，则 $x\notin F_n$，故 $x\in E_{n,k}$ 对唯一的 $k$ 成立，且
\begin{align*}
0\leqslant f(x)-s_n(x)<\frac{k}{2^n}-\frac{k-1}{2^n}=\frac1{2^n}\to0\qquad(n\to\infty).
\end{align*}
故 $s_n(x)\to f(x)$。

有界时的一致收敛。设 $0\leqslant f\leqslant M$。对 $n>M$ 有 $F_n=\varnothing$，故对每个 $x\in X$ 有 $0\leqslant f(x)-s_n(x)<2^{-n}$，于是 $\sup_{x\in X}|f(x)-s_n(x)|\leqslant2^{-n}\to0$。

一般情形。写 $f=f^+-f^-$，其中 $f^\pm\triangleq\max\{\pm f,0\}$ 非负可测（\cref{corollary:ch02-可测运算}(i)）。由上述非负情形，存在非降简单函数列 $s_n^+\uparrow f^+$ 与 $s_n^-\uparrow f^-$。令 $s_n\triangleq s_n^+-s_n^-$；在公共划分上作差可知 $s_n$ 为简单函数，且 $s_n\to f^+ -f^-=f$ 逐点。若 $f$ 有界，则 $f^\pm\leqslant\|f\|_\infty$，故 $s_n^\pm\to f^\pm$ 一致，从而 $s_n\to f$ 一致。此时 $0\leqslant s_n\leqslant f$ 一般不成立，(i) 只对 $f\geqslant0$ 断言。
\end{proof}

% SUPPLEMENT: parent=theorem:ch02-简单逼近; kind=motivation; 说明简单函数逼近的用途、非降性的意义、有界性对一致收敛的必要性、构造对水平集的依赖以及几乎处处定义情形的边界
% SUPPLEMENT: parent=theorem:ch02-简单函数逼近; kind=motivation; 别名标签，与 theorem:ch02-简单逼近 共用下一条注记
\begin{remark}\phantomsection\label{remark:ch02-简单逼近}
\begin{enumerate}[(i)]
\item 引入动机。本定理把一般可测函数化为一列简单函数的逐点极限。非负情形还给出非降性与 $0\leqslant s_n\leqslant f$，这使积分序列 $\int_X s_n\,\mathrm{d}\mu$ 单调收敛于 $\int_X f\,\mathrm{d}\mu$，因而是定义一般可测函数积分与证明单调收敛定理的基础。
\item 各款的作用。(i) 的非降性保证积分的单调性可由简单函数情形传递到极限；(ii) 的一致收敛说明简单函数在 $L^\infty(X)$ 中稠密。有界性不可省：对 $f(x)=x$（$X=\mathbb{R}$），上述构造在 $F_n=\{f\geqslant n\}$ 上取 $s_n=n$，故 $\sup_X|f-s_n|=\infty$，一致收敛不成立。
\item 构造的依赖性。$\{s_n\}$ 依赖于 $f$ 的水平集 $\{f\geqslant k/2^n\}$；这些集合的可测性由\cref{proposition:ch02-可测判据}保证。若 $f$ 只在 $N^c$ 上定义（$N$ 为零测集），则同样的构造在 $N^c$ 上逐点收敛，结论按\cref{definition:ch02-a.e.可测函数}理解。
\item 适用边界。断言的是逐点收敛，不是依测度收敛；一般情形（$f$ 变号）不能保证 $s_n$ 非降；由 $0\leqslant s_n^\pm\leqslant f^\pm$ 仍可得 $|s_n|\leqslant|f|$。若 $f$ 取扩充实值 $\pm\infty$，则需在 $\{|f|=\infty\}$ 上另行处理，简单函数逼近在有限值部分成立。
\end{enumerate}
\end{remark}

为后续需要，考虑定义在可测空间乘积上或取值于可测空间乘积的函数的可测性（乘积 $\sigma$-代数的定义见\cref{definition:ch01-乘积sigma代数}）。

\begin{definition}\label{definition:ch02-投影分量迹}
给定集合 $X$、$X_1$、$X_2$。
\begin{enumerate}[(i)]
\item 映射 $p_k:X_1\times X_2\to X_k$，$p_k(x_1,x_2)\triangleq x_k$，称为 $X_1\times X_2$ 到 $X_k$ 的投影（$k=1,2$）。
\item 设 $f:X\to X_1\times X_2$。映射 $f_k:X\to X_k$，$f_k\triangleq p_k\circ f$，称为 $f$ 的第 $k$ 分量（$k=1,2$）。
\item 设 $f:X_1\times X_2\to X$。对任意 $x_1\in X_1$，映射 $f_{x_1}:X_2\to X$，$f_{x_1}(x_2)\triangleq f(x_1,x_2)$，称为 $f$ 的 $x_1$-迹。对任意 $x_2\in X_2$，映射 $f_{x_2}:X_1\to X$，$f_{x_2}(x_1)\triangleq f(x_1,x_2)$，称为 $f$ 的 $x_2$-迹。
\end{enumerate}
\end{definition}

% SUPPLEMENT: parent=definition:ch02-投影分量迹; kind=motivation; 说明投影、分量与迹的引入动机、各对象的作用、示性函数的迹与截口的关系，以及无穷乘积情形与联合可测性的边界
\begin{remark}\phantomsection\label{remark:ch02-特征函数截面}
\phantomsection\label{remark:ch02-截口}
\begin{enumerate}[(i)]
\item 定义的对象。(i) 中的 $p_k$ 只依赖于乘积空间的集合结构；(ii) 中的分量 $f_k$ 把取值于乘积空间的映射分解为有限多个标量型映射；(iii) 中的迹把定义在乘积空间上的函数对一个变量固定、对另一个变量求值。
\item 示性函数的迹。设 $\chi_E$ 为 $E\subseteq X_1\times X_2$ 的特征函数。对任意 $x_1\in X_1$ 与 $x_2\in X_2$，$(\chi_E)_{x_1}(x_2)=\chi_E(x_1,x_2)$ 等于 $1$ 当且仅当 $(x_1,x_2)\in E$，即 $x_2\in E_{x_1}\triangleq\{x_2\in X_2\mid(x_1,x_2)\in E\}$；故 $(\chi_E)_{x_1}=\chi_{E_{x_1}}$。同理，对任意 $x_2\in X_2$，$(\chi_E)_{x_2}=\chi_{E_{x_2}}$，其中 $E_{x_2}\triangleq\{x_1\in X_1\mid(x_1,x_2)\in E\}$。
\item 引入动机。分量与迹把乘积空间上的可测性问题化归为较低维空间上的问题：\cref{proposition:ch02-乘积可测分量}用分量刻画取值于乘积空间的可测性，\cref{proposition:ch02-迹可测}用迹刻画定义在乘积空间上的可测性，二者是 Tonelli 定理中"截面可测"与"累次积分"的技术准备。
\item 适用边界。分量与迹的定义只用到集合的Descartes积，不需要拓扑或测度。对无穷乘积，分量可测一般推不出联合可测（\cref{proposition:ch02-乘积可测分量}只对有限乘积成立）；反之，各迹可测也推不出联合可测。
\end{enumerate}
\end{remark}

\begin{lemma}\label{lemma:ch02-投影可测}
设 $(X_1\times X_2,\mathscr A_1\times\mathscr A_2)$ 为可测空间 $(X_1,\mathscr A_1)$ 与 $(X_2,\mathscr A_2)$ 的乘积。则对每个 $k=1,2$，投影 $p_k$ 为 $(\mathscr A_1\times\mathscr A_2,\mathscr A_k)$-可测。
\end{lemma}

\begin{proof}
对任意 $E_1\in\mathscr A_1$，
\begin{align*}
p_1^{-1}(E_1)\triangleq \{(x_1,x_2)\in X_1\times X_2\mid x_1\in E_1\}=E_1\times X_2,
\end{align*}
即可测矩形，故 $p_1^{-1}(E_1)\in\mathscr A_1\times\mathscr A_2$。由\cref{lemma:ch02-生成集可测}（取 $\mathscr C'=\mathscr A_1$，$\sigma_0(\mathscr A_1)=\mathscr A_1$）得 $p_1$ 为 $(\mathscr A_1\times\mathscr A_2,\mathscr A_1)$-可测。对 $p_2$ 的论证相同：$p_2^{-1}(E_2)=X_1\times E_2$ 为可测矩形。
\end{proof}

% SUPPLEMENT: parent=lemma:ch02-投影可测; kind=motivation; 说明投影可测性的用途、乘积 $\sigma$-代数由可测矩形生成这一结构的作用，以及投影可测不等于投影保持可测性的边界
\begin{remark}\phantomsection\label{remark:ch02-投影可测}
\begin{enumerate}[(i)]
\item 引入动机。投影可测是乘积 $\sigma$-代数的基本性质。由\cref{proposition:ch02-复合可测}，它把"$f$ 的每个分量可测"与"$f$ 关于乘积 $\sigma$-代数可测"联系起来（\cref{proposition:ch02-乘积可测分量}）。
\item 条件的作用。证明只用到 $\mathscr A_1\times\mathscr A_2$ 含全部可测矩形 $E_1\times X_2$、$X_1\times E_2$。若把乘积空间赋予比 $\mathscr A_1\times\mathscr A_2$ 更小的 $\sigma$-代数，则投影未必可测。
\item 边界情形。$p_1$ 可测并不意味着 $p_1$ 把可测集映为可测集：投影定理（Suslin 理论）指出 Borel 集的连续像未必是 Borel 集。可测性只关乎原像，这正是\cref{definition:ch02-可测函数}的要点。
\end{enumerate}
\end{remark}

\begin{proposition}\label{proposition:ch02-乘积可测分量}
设 $(X,\mathscr A)$、$(X_1,\mathscr A_1)$、$(X_2,\mathscr A_2)$ 为可测空间，$(X_1\times X_2,\mathscr A_1\times\mathscr A_2)$ 为 $(X_1,\mathscr A_1)$ 与 $(X_2,\mathscr A_2)$ 的乘积。设 $f:X\to X_1\times X_2$。则 $f$ 为 $(\mathscr A,\mathscr A_1\times\mathscr A_2)$-可测当且仅当每个分量 $f_k$ 为 $(\mathscr A,\mathscr A_k)$-可测（$k=1,2$）。
\end{proposition}

\begin{proof}
若 $f$ 为 $(\mathscr A,\mathscr A_1\times\mathscr A_2)$-可测，则由\cref{proposition:ch02-复合可测}与\cref{lemma:ch02-投影可测}知每个 $f_k=p_k\circ f$ 为 $(\mathscr A,\mathscr A_k)$-可测。

反之，设 $f_k$ 为 $(\mathscr A,\mathscr A_k)$-可测，从而对任意 $E_k\in\mathscr A_k$ 有 $f_k^{-1}(E_k)\in\mathscr A$。令 $\mathscr R$ 为可测矩形族，即 $\mathscr R\triangleq\{E_1\times E_2\mid E_k\in\mathscr A_k\ (k=1,2)\}$。先证对任意 $E\in\mathscr R$ 有 $f^{-1}(E)\in\mathscr A$。设 $E=E_1\times E_2$，则
\begin{align*}
f^{-1}(E)&=\{x\in X\mid(f_1(x),f_2(x))\in E\}\\
&=\{x\in X\mid f_1(x)\in E_1\}\cap\{x\in X\mid f_2(x)\in E_2\}
=f_1^{-1}(E_1)\cap f_2^{-1}(E_2)\in\mathscr A.
\end{align*}
因 $\mathscr A_1\times\mathscr A_2\triangleq\sigma_0(\mathscr R)$，由\cref{lemma:ch02-生成集可测}即得 $f$ 为 $(\mathscr A,\mathscr A_1\times\mathscr A_2)$-可测。
\end{proof}

% SUPPLEMENT: parent=proposition:ch02-乘积可测分量; kind=motivation; 说明分量刻画的用途、乘积 $\sigma$-代数由可测矩形生成的作用，以及无穷乘积中分量可测不充分的反例
\begin{remark}\phantomsection\label{remark:ch02-分量可测}
\begin{enumerate}[(i)]
\item 引入动机。本命题把取值于乘积空间的映射的可测性化归为有限多个分量的可测性。它是\cref{proposition:ch02-迹可测}的证明工具，也是把向量值可测函数逐分量处理的基础。
\item 条件的作用。充分性方向用到 $\mathscr A_1\times\mathscr A_2$ 由可测矩形生成，因而由\cref{lemma:ch02-生成集可测}只需验证可测矩形的原像；必要性方向用到投影可测与复合封闭。两个方向缺一不可：没有乘积 $\sigma$-代数的生成结构，充分性不成立；没有投影可测，必要性不成立。
\item 边界情形。对无穷乘积 $\prod_{n\in\mathbb{N}}X_n$ 及其乘积 $\sigma$-代数，每个分量可测不能推出联合可测，因为乘积 $\sigma$-代数由有限维柱集生成，而一般集合可能不是有限维柱集的可数极限。本命题只对两个因子陈述，多个因子的情形由归纳得到。
\end{enumerate}
\end{remark}

\begin{proposition}\label{proposition:ch02-迹可测}
设 $(X_1,\mathscr A_1)$、$(X_2,\mathscr A_2)$、$(X,\mathscr A)$ 为可测空间，$f:X_1\times X_2\to X$ 为 $(\mathscr A_1\times\mathscr A_2,\mathscr A)$-可测。则：
\begin{enumerate}[(i)]
\item 对任意 $x_1\in X_1$，$f$ 的 $x_1$-迹为 $(\mathscr A_2,\mathscr A)$-可测；
\item 对任意 $x_2\in X_2$，$f$ 的 $x_2$-迹为 $(\mathscr A_1,\mathscr A)$-可测。
\end{enumerate}
\end{proposition}

% SUPPLEMENT: parent=proposition:ch02-迹可测; kind=proof; 原书对 (ii) 称"论证类似"而未写出，此处补出两个方向的嵌入映射、常值映射的可测性与复合论证
\begin{proof}
(i) 固定 $x_1\in X_1$，令 $g:X_2\to X_1\times X_2$，$g(x_2)\triangleq(x_1,x_2)$。则 $f_{x_1}=f\circ g$。$g$ 的两个分量分别为常值映射 $g_1\equiv x_1$ 与恒等映射 $g_2=\mathrm{id}_{X_2}$：前者满足 $g_1^{-1}(E_1)\in\{\varnothing,X_2\}\subseteq\mathscr A_2$ 对任意 $E_1\in\mathscr A_1$，故为 $(\mathscr A_2,\mathscr A_1)$-可测；后者为 $(\mathscr A_2,\mathscr A_2)$-可测。由\cref{proposition:ch02-乘积可测分量}得 $g$ 为 $(\mathscr A_2,\mathscr A_1\times\mathscr A_2)$-可测，再由\cref{proposition:ch02-复合可测}得 $f_{x_1}$ 为 $(\mathscr A_2,\mathscr A)$-可测。

(ii) 固定 $x_2\in X_2$，令 $h:X_1\to X_1\times X_2$，$h(x_1)\triangleq(x_1,x_2)$。则 $f_{x_2}=f\circ h$。$h$ 的两个分量分别为 $\mathrm{id}_{X_1}$ 与常值映射 $x_1\mapsto x_2$，前者为 $(\mathscr A_1,\mathscr A_1)$-可测，后者为 $(\mathscr A_1,\mathscr A_2)$-可测。同上，由\cref{proposition:ch02-乘积可测分量}与\cref{proposition:ch02-复合可测}得 $f_{x_2}$ 为 $(\mathscr A_1,\mathscr A)$-可测。
\end{proof}

% SUPPLEMENT: parent=proposition:ch02-迹可测; kind=motivation; 说明迹可测性的用途、关于乘积 $\sigma$-代数的假设的作用，以及逆命题不成立的边界
\begin{remark}\phantomsection\label{remark:ch02-截面可测}
\phantomsection\label{remark:ch02-截口可测}
\begin{enumerate}[(i)]
\item 引入动机。本命题说明联合可测函数的所有迹都单变量可测，这是 Tonelli 定理中"对几乎每个 $x_1$，函数 $x_2\mapsto f(x_1,x_2)$ 可测"的依据，也是 Fubini 定理中累次积分有意义的前提。
\item 示性函数的情形。设 $\chi_E$ 为 $E\subseteq X_1\times X_2$ 的特征函数。若 $E\in\mathscr A_1\times\mathscr A_2$，则 $\chi_E$ 为 $\mathscr A_1\times\mathscr A_2$-可测，由 (i) 知 $(\chi_E)_{x_1}=\chi_{E_{x_1}}$ 为 $\mathscr A_2$-可测（此处用到 $E_{x_1}\in\mathscr A_2$，见\cref{proposition:ch01-截面}）。同理，对任意 $x_2\in X_2$，$(\chi_E)_{x_2}=\chi_{E_{x_2}}$ 为 $\mathscr A_1$-可测。
\item 逆命题不成立。各迹可测一般推不出联合可测：在连续统假设下存在 Sierpiński 集 $E\subseteq\mathbb{R}^2$，其每个截面至多可数（从而各迹可测），但 $E$ 不是 Lebesgue 可测集，故 $\chi_E$ 不可测。这说明本命题不能反向使用。
\end{enumerate}
\end{remark}

\subsection{收敛结果}\label{subsection:ch02-收敛结果}

设 $(X,\mathscr A,\mu)$ 为测度空间，$f,f_n:X\to\mathbb{R}$ 为 $\mathscr A$-可测函数（$n\in\mathbb{N}$）。由\cref{corollary:ch02-可测运算}，水平集 $\{|f_n-f|>\varepsilon\}$ 属于 $\mathscr A$，故其测度有意义。下列收敛概念（统一记为 $f_n\to f$）是熟知的：
\begin{itemize}
\item $f_n$ \textbf{依测度}收敛于 $f$，若对所有 $\varepsilon>0$ 有 $\lim_{n\to\infty}\mu(\{x\in X\mid|f_n(x)-f(x)|>\varepsilon\})=0$；
\item $f_n$ \textbf{几乎一致}收敛于 $f$，若对任意 $\delta>0$，存在 $E\in\mathscr A$ 使 $\mu(E^c)<\delta$ 且
\begin{align}
\lim_{n\to\infty}\sup_{x\in E}|f_n(x)-f(x)|=0. \label{align:ch02-2.3}
\end{align}
\end{itemize}

回顾上述概念之间的关系。

\begin{proposition}\label{proposition:ch02-收敛关系}
设 $(X,\mathscr A,\mu)$ 为测度空间，$f,f_n:X\to\mathbb{R}$ 为 $\mathscr A$-可测函数（$n\in\mathbb{N}$）。
\begin{enumerate}[(i)]
\item 设 $\mu(X)<\infty$。若 $f_n\to f$ $\mu$-a.e.，则 $f_n\to f$ 依测度。
\item 若 $f_n\to f$ 依测度，则存在子列 $\{f_{n_k}\}\subseteq\{f_n\}$ 使 $f_{n_k}\to f$ $\mu$-a.e.。
\item 若 $f_n\to f$ 几乎一致，则 $f_n\to f$ $\mu$-a.e. 且依测度。
\end{enumerate}
\end{proposition}

% SUPPLEMENT: parent=proposition:ch02-收敛关系; kind=proof; 原书仅回顾结论而未给证明，此处补出有限测度下的上连续性论证、依测度收敛取子列用可数次可加性、以及几乎一致收敛推出几乎处处与依测度收敛的完整证明
\begin{proof}
(i) 设 $\mu(X)<\infty$ 且 $f_n\to f$ $\mu$-a.e.。取零测集 $N\in\mathscr A$ 使 $\lim_{n\to\infty}|f_n(x)-f(x)|=0$ 对每个 $x\in N^c$ 成立。固定 $\varepsilon>0$，令 $A_n\triangleq\{|f_n-f|>\varepsilon\}\in\mathscr A$。若 $x\in\limsup_n A_n$，则 $x$ 属于无穷多个 $A_n$，即 $|f_n(x)-f(x)|>\varepsilon$ 对无穷多个 $n$ 成立，故 $x\notin N^c$，即 $\limsup_nA_n\subseteq N$。于是
\begin{align*}
\mu\Bigl(\limsup_{n\to\infty}A_n\Bigr)=0.
\end{align*}
另一方面，$\{A_n\}$ 的并满足 $\bigcup_{k\geqslant n}A_k\supseteq A_n$，故 $\mu(A_n)\leqslant\mu\bigl(\bigcup_{k\geqslant n}A_k\bigr)$，从而
\begin{align*}
\limsup_{n\to\infty}\mu(A_n)\leqslant\lim_{n\to\infty}\mu\Bigl(\bigcup_{k\geqslant n}A_k\Bigr)=\mu\Bigl(\bigcap_{n=1}^\infty\bigcup_{k\geqslant n}A_k\Bigr)=\mu\Bigl(\limsup_{n\to\infty}A_n\Bigr)=0,
\end{align*}
这里第二个等号用到 $\mu\bigl(\bigcup_{k\geqslant1}A_k\bigr)\leqslant\mu(X)<\infty$ 与测度的上连续性。故 $\mu(A_n)\to0$，即 $f_n\to f$ 依测度。

(ii) 设 $f_n\to f$ 依测度。对每个 $k\in\mathbb{N}$，由依测度收敛，存在 $m_k\in\mathbb{N}$ 使 $n\geqslant m_k$ 时
\begin{align*}
\mu\bigl(\{|f_n-f|\geqslant2^{-k}\}\bigr)<2^{-k}.
\end{align*}
归纳地取 $n_1\triangleq m_1$、$n_{k+1}\triangleq\max\{m_{k+1},n_k+1\}$，则 $\{n_k\}$ 严格增，故 $\{f_{n_k}\}$ 为子列，且 $B_k\triangleq\{|f_{n_k}-f|\geqslant2^{-k}\}$ 满足 $\mu(B_k)<2^{-k}$。由可数次可加性，对任意 $m\in\mathbb{N}$ 有
\begin{align*}
\mu\Bigl(\bigcup_{k=m}^\infty B_k\Bigr)\leqslant\sum_{k=m}^\infty\mu(B_k)<\sum_{k=m}^\infty2^{-k}=2^{1-m},
\end{align*}
故 $\mu\bigl(\limsup_{k\to\infty}B_k\bigr)\leqslant2^{1-m}$ 对一切 $m$ 成立，于是 $\mu\bigl(\limsup_kB_k\bigr)=0$。若 $x\notin\limsup_kB_k$，则 $x$ 只属于有限多个 $B_k$，故存在 $k_0$ 使 $k\geqslant k_0$ 时 $|f_{n_k}(x)-f(x)|<2^{-k}$，从而 $f_{n_k}(x)\to f(x)$。故 $f_{n_k}\to f$ $\mu$-a.e.。

(iii) 设 $f_n\to f$ 几乎一致。对每个 $m\in\mathbb{N}$，存在 $E_m\in\mathscr A$ 使 $\mu(E_m^c)<1/m$ 且 $\sup_{x\in E_m}|f_n(x)-f(x)|\to0$。令 $E\triangleq\bigcup_{m=1}^\infty E_m$。由 $E^c\subseteq E_m^c$ 得 $\mu(E^c)\leqslant\mu(E_m^c)<1/m$ 对一切 $m$ 成立，故 $\mu(E^c)=0$；且对每个 $x\in E$，存在 $m$ 使 $x\in E_m$，故 $f_n(x)\to f(x)$。于是 $f_n\to f$ $\mu$-a.e.。

再证依测度收敛。给定 $\varepsilon>0$ 与 $\delta>0$，取 $m$ 使 $1/m<\delta$。由 $E_m$ 上的一致收敛，存在 $\bar n\in\mathbb{N}$ 使 $n\geqslant\bar n$ 时 $\sup_{x\in E_m}|f_n(x)-f(x)|\leqslant\varepsilon$，即 $E_m\subseteq\{|f_n-f|\leqslant\varepsilon\}$。于是 $\{|f_n-f|>\varepsilon\}\subseteq E_m^c$，从而
\begin{align*}
\mu\bigl(\{|f_n-f|>\varepsilon\}\bigr)\leqslant\mu(E_m^c)<1/m<\delta.
\end{align*}
由 $\varepsilon$ 与 $\delta$ 的任意性，$f_n\to f$ 依测度。
\end{proof}

% SUPPLEMENT: parent=proposition:ch02-收敛关系; kind=motivation; 说明三种收敛的关系、(i) 中有限测度不可省的反例、(ii) 中子列不可省的反例，以及 Egorov 定理的定位
\begin{remark}\phantomsection\label{remark:ch02-收敛关系}
\begin{enumerate}[(i)]
\item 引入动机。依测度收敛、几乎一致收敛与几乎处处收敛是积分理论中三种基本收敛方式。本命题给出 三种收敛之间的部分蕴含关系；Egorov 定理（\cref{theorem:ch02-Egorov}）则在有限测度空间上把 (i) 反向加强为几乎一致收敛。
\item (i) 的有限测度不可省。取 $(X,\mathscr A,\mu)=(\mathbb{R},\mathscr B(\mathbb{R}),\lambda)$ 与 $f_n\triangleq\chi_{[n,n+1]}$、$f\equiv0$。则对每个 $x\in\mathbb{R}$，当 $n>x$ 时 $f_n(x)=0$，故 $f_n\to0$ 处处成立；但 $\lambda(\{|f_n|>1/2\})=\lambda([n,n+1])=1$ 不趋于 $0$，故不依测度收敛。论证中上连续性用到 $\mu\bigl(\bigcup_{k\geqslant1}A_k\bigr)<\infty$，有限性正是此处所需。
\item (ii) 的子列不可省。在 $(0,1)$ 上把区间 $[0,1)$ 按 $N=1,2,\cdots$ 分成 $N$ 个等长区间 $I_{N,j}\triangleq[(j-1)/N,j/N)$（$j=1,\cdots,N$），并按 $N$ 递增、同一 $N$ 内 $j$ 递增的次序排成序列 $\{I_n\}$，令 $f_n\triangleq\chi_{I_n}$。则 $\lambda(I_n)=1/N\to0$，故 $f_n\to0$ 依测度；但每个 $x\in(0,1)$ 在每个尺度 $N$ 下恰属于一个 $I_{N,j}$，故 $f_n(x)=1$ 对无穷多个 $n$ 成立，$\{f_n(x)\}$ 不收敛。因此依测度收敛只能保证存在几乎处处收敛的子列。
\item (iii) 的强度。几乎一致收敛强于几乎处处收敛：前者要求去掉任意小测度集后一致收敛。反过来，几乎处处收敛在有限测度空间上蕴含几乎一致收敛（Egorov），但在无限测度空间上不成立（\cref{theorem:ch02-Egorov}及其后的注记）。
\end{enumerate}
\end{remark}

下列结果在有限测度空间上部分逆转\cref{proposition:ch02-收敛关系}(iii)的蕴含方向。

\begin{theorem}[Egorov]\label{theorem:ch02-Egorov}
设 $(X,\mathscr A,\mu)$ 为有限测度空间，即 $\mu(X)<\infty$，并设 $f,f_n:X\to\mathbb{R}$ 为 $\mathscr A$-可测函数（$n\in\mathbb{N}$）。若 $f_n\to f$ $\mu$-a.e.，则 $f_n\to f$ 几乎一致。
\end{theorem}

% SUPPLEMENT: parent=theorem:ch02-Egorov; kind=proof; 原证明中"容易看出 $G_k\subseteq N$"与上连续性、$\delta/2^{k}$ 求和等步骤被压缩，此处补出逐项验证并写清一致收敛估计
\begin{proof}
由假设存在 $\mu$-零测集 $N\in\mathscr A$，使对每个 $x\in N^c$ 有 $\lim_{n\to\infty}|f_n(x)-f(x)|=0$。对 $k,n,p\in\mathbb{N}$，令
\begin{align*}
E_{n,k}\triangleq \{x\in X\mid|f_n(x)-f(x)|\geqslant 1/k\}\in\mathscr A,\qquad
F_{p,k}\triangleq \bigcup_{n=p}^\infty E_{n,k}\in\mathscr A,
\end{align*}
其可测性由\cref{corollary:ch02-可测运算}得到。

第一步：对每个 $k\in\mathbb{N}$，集合 $G_k\triangleq\bigcap_{p=1}^\infty F_{p,k}$ 含于 $N$。固定 $k\in\mathbb{N}$ 与 $x\in G_k$。对每个 $p\in\mathbb{N}$，由 $x\in F_{p,k}$ 与 $F_{p,k}$ 的定义，存在 $n_p\geqslant p$ 使 $x\in E_{n_p,k}$，即 $|f_{n_p}(x)-f(x)|\geqslant1/k$。于是 $n_p\to\infty$ 而 $|f_{n_p}(x)-f(x)|\geqslant1/k$，故 $|f_n(x)-f(x)|$ 不收敛于 $0$，即 $x\notin N^c$。故 $x\in N$。

第二步：对每个 $k\in\mathbb{N}$ 有 $\mu(G_k)=0$。事实上，$G_k$ 可测且 $G_k\subseteq N$，而 $\mu(N)=0$。

第三步：对固定 $k\in\mathbb{N}$，序列 $\{F_{p,k}\}_{p\in\mathbb{N}}$ 非增，且
\begin{align*}
\mu(F_{1,k})\leqslant\mu(X)<\infty.
\end{align*}
由测度的上连续性得
\begin{align*}
\lim_{p\to\infty}\mu(F_{p,k})=\mu\Bigl(\bigcap_{p=1}^\infty F_{p,k}\Bigr)=\mu(G_k)=0.
\end{align*}

第四步：给定 $\delta>0$。对每个 $k\in\mathbb{N}$，由第三步存在 $p_k\in\mathbb{N}$ 使
\begin{align*}
\mu(F_{p_k,k})<\frac{\delta}{2^{k+1}}.
\end{align*}
令 $F_\delta\triangleq\bigcup_{k=1}^\infty F_{p_k,k}$。由可数次可加性，
\begin{align*}
\mu(F_\delta)\leqslant\sum_{k=1}^\infty\mu(F_{p_k,k})<\sum_{k=1}^\infty\frac{\delta}{2^{k+1}}=\frac{\delta}{2}<\delta.
\end{align*}
令 $E\triangleq F_\delta^c\in\mathscr A$，则 $\mu(E^c)=\mu(F_\delta)<\delta$。

第五步：$f_n\to f$ 在 $E$ 上一致。由
\begin{align*}
E=F_\delta^c=\bigcap_{k=1}^\infty F_{p_k,k}^c=\bigcap_{k=1}^\infty\bigcap_{n=p_k}^\infty E_{n,k}^c
\end{align*}
知：固定 $x\in E$ 与 $k\in\mathbb{N}$，则对所有 $n\geqslant p_k$ 有 $x\notin E_{n,k}$，即 $|f_n(x)-f(x)|<1/k$。因此对每个 $k\in\mathbb{N}$ 与所有 $n\geqslant p_k$ 有
\begin{align*}
\sup_{x\in E}|f_n(x)-f(x)|\leqslant\frac1k,
\end{align*}
从而
\begin{align*}
\limsup_{n\to\infty}\sup_{x\in E}|f_n(x)-f(x)|\leqslant\frac1k.
\end{align*}
由 $k$ 的任意性得该上极限为 $0$，即 $\lim_{n\to\infty}\sup_{x\in E}|f_n(x)-f(x)|=0$。结合 $\mu(E^c)<\delta$ 与 $\delta$ 的任意性，$f_n\to f$ 几乎一致。
\end{proof}

% SUPPLEMENT: parent=theorem:ch02-Egorov; kind=motivation; 说明 Egorov 定理的用途、有限测度不可省的反例、可测性假设的作用、不能加强为全空间一致收敛的边界以及 $\delta$ 与 $k$ 选取的依赖关系
\begin{remark}\phantomsection\label{remark:ch02-Egorov}
\begin{enumerate}[(i)]
\item 引入动机。Egorov 定理把几乎处处收敛加强为几乎一致收敛：除去一个测度任意小的集合后，收敛是一致的。它使"几乎处处收敛"在有限测度空间上可以按几乎一致收敛处理，是控制收敛定理与 $L^p$ 收敛论证中提取一致收敛子列的工具。
\item 有限测度不可省。取 $(X,\mathscr A,\mu)=(\mathbb{R},\mathscr B(\mathbb{R}),\lambda)$ 与 $f_n\triangleq\chi_{[n,n+1]}$、$f\equiv0$。则 $f_n\to0$ 处处成立。若 $f_n\to0$ 几乎一致，则存在 $E\in\mathscr B(\mathbb{R})$ 使 $\lambda(E^c)<\delta$ 且 $\sup_{x\in E}|f_n(x)|\to0$。于是存在 $\bar n$ 使 $n\geqslant\bar n$ 时 $E\cap[n,n+1]=\varnothing$，即 $[n,n+1]\subseteq E^c$ 对一切 $n\geqslant\bar n$ 成立，从而 $\lambda(E^c)=\infty$，矛盾。故结论在无限测度空间上不成立。证明中第三步的上连续性用到 $\mu(F_{1,k})\leqslant\mu(X)<\infty$，这正是有限性被使用之处。
\item 可测性假设的作用。第一步中 $E_{n,k}$、$F_{p,k}$、$G_k$ 的可测性来自 $f_n,f$ 的 $\mathscr A$-可测性（\cref{corollary:ch02-可测运算}）。若函数不可测，则这些集合未必属于 $\mathscr A$，"几乎一致收敛"的定义本身也无法陈述。
\item 不能加强为全空间一致收敛。取 $X=(0,1)$、$\mu=\lambda$、$f_n\triangleq\chi_{(0,1/n)}$、$f\equiv0$。则 $f_n\to0$ 处处成立，但 $\sup_{x\in(0,1)}|f_n(x)|=1$ 对一切 $n$ 成立，故不一致收敛。Egorov 定理只断言可以去掉任意小测度的集合。
\item 量词的依赖关系。结论中的 $E$ 依赖于 $\delta$，$p_k$ 依赖于 $k$ 与 $\delta$，不能取到与 $\delta$ 无关的 $E$；第五步中的估计 $\sup_E|f_n-f|\leqslant1/k$ 对 $n\geqslant p_k$ 成立，$p_k$ 随 $k$ 增大而增大。
\end{enumerate}
\end{remark}

若以容量代替测度，类似概念与结果成立。设 $(X,d)$ 为度量空间，$\mathscr F\subseteq\mathscr P(X)$ 如\cref{definition:ch01-容量}所述，$C:\mathscr F\to[0,\infty]$ 为容量。以下约定 $\mathscr F$ 对补集封闭，即 $E\in\mathscr F\Rightarrow X\setminus E\in\mathscr F$（当 $\mathscr F=\mathscr P(X)$ 时自动成立）；否则定义中的 $C(E^c)$ 无意义。由\hyperref[remark:ch01-容量与外测度]{关于容量的注记}，开集族含于 $\mathscr F$，故对任意 $\varepsilon>0$，水平集 $\{|f_n-f|>\varepsilon\}$ 属于 $\mathscr F$，$C$ 在其上有定义。

\begin{definition}\label{definition:ch02-容量收敛}
设 $(X,d)$ 为度量空间，$\mathscr F\subseteq\mathscr P(X)$ 如\cref{definition:ch01-容量}所述且对补集封闭，$C:\mathscr F\to[0,\infty]$ 为容量。设 $f,f_n:X\to\mathbb{R}$（$n\in\mathbb{N}$）。称：
\begin{enumerate}[(i)]
\item $f_n$ 在 $X$ 上 \textbf{$C$-拟处处}收敛于 $f$，若存在 $N\in\mathscr F$ 使 $C(N)=0$ 且对任意 $x\in X\setminus N$ 有 $f_n(x)\to f(x)$；
\item $f_n$ \textbf{依容量}收敛于 $f$，若对所有 $\varepsilon>0$ 有 $\lim_{n\to\infty}C(\{|f_n-f|>\varepsilon\})=0$；
\item $f_n$ \textbf{$C$-拟一致}收敛于 $f$，若对任意 $\delta>0$，存在 $E\in\mathscr F$ 使 $C(E^c)<\delta$ 且\eqref{align:ch02-2.3}成立。
\end{enumerate}
\end{definition}

% SUPPLEMENT: parent=definition:ch02-容量收敛; kind=motivation; 说明容量收敛的三种概念的动机、$\mathscr F$ 对补集封闭这一约定的必要性、开集属于 $\mathscr F$ 的良定义性，以及与测度情形的对照
\begin{remark}\phantomsection\label{remark:ch02-依容量判据}
\begin{enumerate}[(i)]
\item 等价刻画。$f_n\to f$ 依容量当且仅当：对任意 $\delta>0$ 与 $\varepsilon>0$，存在 $\bar n\in\mathbb{N}$，使对所有 $n\geqslant\bar n$ 有 $C(\{|f_n-f|\geqslant\varepsilon\})<\delta$。事实上，若依容量收敛，则给定 $\delta,\varepsilon$，由 $\{|f_n-f|\geqslant\varepsilon\}\subseteq\{|f_n-f|>\varepsilon/2\}$ 与 $C$ 的单调性，存在 $\bar n$ 使 $n\geqslant\bar n$ 时
\begin{align*}
C(\{|f_n-f|\geqslant\varepsilon\})\leqslant C(\{|f_n-f|>\varepsilon/2\})<\delta.
\end{align*}
反之，若上述 $\delta$-$\varepsilon$ 条件成立，则给定 $\varepsilon>0$ 与 $\delta>0$，取条件中的 $\varepsilon$ 与 $\delta$，得 $\bar n$ 使 $n\geqslant\bar n$ 时 $C(\{|f_n-f|>\varepsilon\})\leqslant C(\{|f_n-f|\geqslant\varepsilon\})<\delta$，故 $C(\{|f_n-f|>\varepsilon\})\to0$。
\item 引入动机。三种 $C$-收敛是测度情形三种收敛的容量版本，用于处理不具备测度结构、但具备容量估计的函数列（例如 Sobolev 空间中按容量收敛的逼近列）。
\item 约定的作用。定义中的 $C(E^c)$ 要求 $E^c\in\mathscr F$，故必须约定 $\mathscr F$ 对补集封闭。水平集 $\{|f_n-f|>\varepsilon\}$ 是开集，由\cref{definition:ch02-容量收敛}前的约定属于 $\mathscr F$；这一点保证"依容量收敛"的定义是良定的。容量只要求单调性与可数次可加性，不要求可加性，故不能照搬测度论中的补集论证与差集运算。
\item 边界情形。$C$ 一般不具有测度的上连续性：对非增集合序列 $\{E_k\}\subseteq\mathscr F$，$C(\bigcap_kE_k)$ 未必等于 $\lim_kC(E_k)$（Choquet 性质只对非增紧集序列成立，见\cref{proposition:ch01-外容量上连续}）。因此\cref{proposition:ch02-容量收敛关系}的证明不能直接沿用测度情形，须改用 $C$ 的可数次可加性。
\end{enumerate}
\end{remark}

下列命题与\cref{proposition:ch02-收敛关系}类似。

\begin{proposition}\label{proposition:ch02-容量收敛关系}
设 $(X,d)$ 为度量空间，$\mathscr F\subseteq\mathscr P(X)$ 如\cref{definition:ch01-容量}所述且对补集封闭，$C:\mathscr F\to[0,\infty]$ 为容量。设 $f,f_n:X\to\mathbb{R}$（$n\in\mathbb{N}$）。
\begin{enumerate}[(i)]
\item 若 $f_n\to f$ $C$-拟一致，则 $f_n\to f$ 依容量且 $C$-q.e.。
\item 若 $f_n\to f$ 依容量，则存在子列 $\{f_{n_k}\}\subseteq\{f_n\}$ 使 $f_{n_k}\to f$ $C$-拟一致。
\end{enumerate}
\end{proposition}

% SUPPLEMENT: parent=proposition:ch02-容量收敛关系; kind=proof; 原证明中水平集的可测性、$C(A_m)$ 的估计、$E\in\mathscr F$ 的选取与 $C$-q.e. 的并集论证均有跳跃，此处补出完整推导
\begin{proof}
(i) 先证依容量收敛。给定 $\varepsilon>0$ 与 $\delta>0$。由 $C$-拟一致收敛，存在 $E\in\mathscr F$ 使 $C(E^c)<\delta$ 且 $\lim_{n\to\infty}\sup_{x\in E}|f_n(x)-f(x)|=0$。于是存在 $\bar n\in\mathbb{N}$，使 $n\geqslant\bar n$ 时 $\sup_{x\in E}|f_n(x)-f(x)|\leqslant\varepsilon$，即 $E\subseteq\{|f_n-f|\leqslant\varepsilon\}$。取补集得 $\{|f_n-f|>\varepsilon\}\subseteq E^c$。因水平集为开集、$E^c\in\mathscr F$，由 $C$ 的单调性，
\begin{align*}
C(\{|f_n-f|>\varepsilon\})\leqslant C(E^c)<\delta.
\end{align*}
故依容量收敛（见\hyperref[remark:ch02-依容量判据]{前一条注记}(i)）。

再证 $C$-q.e. 收敛。对每个 $k\in\mathbb{N}$，由 $C$-拟一致收敛，存在 $E_k\in\mathscr F$ 使 $C(E_k^c)<1/k$ 且 $\sup_{x\in E_k}|f_n(x)-f(x)|\to0$。令 $E\triangleq\bigcup_{k=1}^\infty E_k\in\mathscr F$。若 $x\in E$，则 $x\in E_k$ 对某个 $k$ 成立，故 $f_n(x)\to f(x)$。又 $E^c\subseteq E_k^c$ 对每个 $k$ 成立，故由单调性 $C(E^c)\leqslant C(E_k^c)<1/k$ 对一切 $k$ 成立，于是 $C(E^c)=0$。故 $f_n\to f$ $C$-q.e.。

(ii) 由\hyperref[remark:ch02-依容量判据]{前述注记}(i)，对每个 $k\in\mathbb{N}$，存在 $m_k\in\mathbb{N}$，使 $n\geqslant m_k$ 时
\begin{align*}
C\bigl(\{|f_n-f|\geqslant2^{-k}\}\bigr)<2^{-k}.
\end{align*}
归纳地取 $n_1\triangleq m_1$、$n_{k+1}\triangleq\max\{m_{k+1},n_k+1\}$，则 $\{n_k\}$ 严格增，故 $\{f_{n_k}\}$ 为子列。令
\begin{align*}
B_k\triangleq\{|f_{n_k}-f|\geqslant2^{-k}\}\in\mathscr F,\qquad
A_m\triangleq\bigcup_{k=m}^\infty B_k\in\mathscr F\quad(m\in\mathbb{N}).
\end{align*}
则 $C(B_k)<2^{-k}$，且由可数次可加性
\begin{align*}
C(A_m)\leqslant\sum_{k=m}^\infty C(B_k)<\sum_{k=m}^\infty2^{-k}=2^{1-m}\to0\qquad(m\to\infty).
\end{align*}
故对任意 $\delta>0$，存在 $\bar m\in\mathbb{N}$ 使 $C(A_{\bar m})<\delta$。令 $E\triangleq A_{\bar m}^c\in\mathscr F$，则 $C(E^c)=C(A_{\bar m})<\delta$。若 $k\geqslant\bar m$ 与 $x\in E$，则 $x\notin B_k$，即 $|f_{n_k}(x)-f(x)|<2^{-k}\leqslant2^{-\bar m}$。于是
\begin{align*}
\sup_{x\in E}|f_{n_k}(x)-f(x)|\leqslant2^{-\bar m}\qquad(k\geqslant\bar m),
\end{align*}
令 $k\to\infty$ 得 $\lim_{k\to\infty}\sup_{x\in E}|f_{n_k}(x)-f(x)|=0$。故 $f_{n_k}\to f$ $C$-拟一致。
\end{proof}

% SUPPLEMENT: parent=proposition:ch02-容量收敛关系; kind=motivation; 说明与测度情形的平行关系、$C$ 的可数次可加性在何处使用、子列不可省以及容量缺少上连续性造成的论证差别
\begin{remark}\phantomsection\label{remark:ch02-容量收敛关系}
\begin{enumerate}[(i)]
\item 引入动机。本命题是容量框架下的基本收敛关系，与\cref{proposition:ch02-收敛关系}平行。它把 $C$-拟一致收敛与依容量收敛、$C$-拟处处收敛联系起来，并在依容量收敛时提取 $C$-拟一致收敛的子列。
\item 条件的作用。证明只用到 $C$ 的单调性与可数次可加性，以及 $\mathscr F$ 对补集封闭（用于 $C(E^c)$ 与 $E=A_{\bar m}^c$ 的选取）。与测度情形不同，这里不能使用上连续性；取而代之的是把可数个"坏集" $B_k$ 并成 $A_m$ 并用可数次可加性控制 $C(A_m)$。
\item 子列不可省。与\cref{proposition:ch02-收敛关系}(ii) 同理，依容量收敛一般不能保证全列 $C$-拟一致收敛；上述证明只对按 $2^{-k}$ 选取的子列给出估计。
\item 边界情形。$C$-拟一致收敛的定义要求 $E\in\mathscr F$ 且 $E^c\in\mathscr F$；若 $\mathscr F$ 不对补集封闭，则这些陈述需要另行解释。此外，结论中的 $E$ 依赖于 $\delta$，$n_k$ 依赖于 $k$。
\end{enumerate}
\end{remark}

\subsection{拟连续函数}\label{subsection:ch02-拟连续}

\begin{theorem}[Lusin]\label{theorem:ch02-Lusin}
设 $X$ 为 Hausdorff 空间，$\mu$ 为 $\mathscr B(X)$ 上的正则测度。设 $f:X\to\mathbb{R}$ 为 $\mathscr B(X)$-可测。则对任意 $E\in\mathscr B(X)$（$\mu(E)<\infty$）与任意 $\delta>0$，存在紧集 $K\subseteq E$ 使 $\mu(E\setminus K)<\delta$，且限制 $f|_K$ 连续。
\end{theorem}

% SUPPLEMENT: parent=theorem:ch02-Lusin; kind=proof; 原证明中"只需考虑 $E=X$"、$A_n\cap C=K_n\cap C$ 的等式与由可数基得到 $f|_C$ 连续的步骤均有压缩，此处补出测度限制的正则性验证、集合等式的双向包含与相对开性的完整论证
\begin{proof}
第一步：化归到有限测度空间。固定 $E\in\mathscr B(X)$ 且 $\mu(E)<\infty$。在 $E$ 上赋迹拓扑与迹 $\sigma$-代数 $\mathscr B(X)\cap E$，令 $\mu_E\triangleq\mu|_{\mathscr B(X)\cap E}$。则 $\mu_E(E)=\mu(E)<\infty$。验证 $\mu_E$ 正则：对 $A\in\mathscr B(X)\cap E$ 与 $\varepsilon>0$，由 $\mu$ 内正则存在紧集 $K\subseteq A$ 使 $\mu(A\setminus K)<\varepsilon$，而 $A\setminus K=(E\setminus K)\cap A$，故 $\mu_E(A\setminus K)<\varepsilon$；由 $\mu$ 外正则存在开集 $U\supseteq A$ 使 $\mu(U\setminus A)<\varepsilon$，而 $U\cap E$ 为 $E$ 中的开集且 $(U\cap E)\setminus A=U\setminus A$，故 $\mu_E((U\cap E)\setminus A)<\varepsilon$。又 $f|_E$ 为 $\mathscr B(X)\cap E$-可测。因此只需证明 $E=X$ 且 $\mu(X)<\infty$ 的情形。

第二步：有限测度情形。设 $\mu(X)<\infty$。取 $\mathbb{R}$ 中开区间的一个可数基 $\{B_n\}_{n\in\mathbb{N}}$。由 $f$ 的 $\mathscr B(X)$-可测性，$f^{-1}(B_n)\in\mathscr B(X)$（$n\in\mathbb{N}$）。固定 $\delta>0$。由 $\mu$ 的正则性，对每个 $n\in\mathbb{N}$，存在紧集 $K_n$ 与开集 $A_n$ 使
\begin{align*}
K_n\subseteq f^{-1}(B_n)\subseteq A_n,\qquad \mu(A_n\setminus K_n)<\frac{\delta}{2^{n+1}}.
\end{align*}
令
\begin{align*}
A\triangleq\bigcup_{n=1}^\infty(A_n\setminus K_n),\qquad C\triangleq A^c.
\end{align*}
则 $A$ 为开集，从而 $C$ 为闭集，且
\begin{align*}
\mu(A)=\mu(C^c)\leqslant\sum_{n=1}^\infty\mu(A_n\setminus K_n)\leqslant\frac{\delta}{2}<\infty.
\end{align*}

第三步：$f|_C$ 连续。对任意 $n\in\mathbb{N}$，由 $A_n\setminus K_n\subseteq A$ 得 $A_n\cap C\subseteq K_n\cap C$；反向包含 $K_n\cap C\subseteq A_n\cap C$ 由 $K_n\subseteq A_n$ 得到。于是
\begin{align*}
A_n\cap C=K_n\cap C\subseteq f^{-1}(B_n)\cap C\subseteq A_n\cap C,
\end{align*}
从而
\begin{align*}
f^{-1}(B_n)\cap C=(f|_C)^{-1}(B_n)=A_n\cap C.
\end{align*}
右端是 $C$ 的相对拓扑中的开集。对任意开集 $V\subseteq\mathbb{R}$，由 $\{B_n\}$ 为基得 $V=\bigcup\{B_n\mid B_n\subseteq V\}$，故
\begin{align*}
(f|_C)^{-1}(V)=\bigcup_{B_n\subseteq V}(f|_C)^{-1}(B_n)
\end{align*}
是 $C$ 的相对拓扑中的开集之并，因而相对开。于是 $f|_C$ 连续。

第四步：收紧为紧集。由 $\mu$ 的内正则性与 $\mu(C)\leqslant\mu(X)<\infty$，存在紧集 $K\subseteq C$ 使 $\mu(C\setminus K)<\delta/2$。于是
\begin{align*}
\mu(X\setminus K)=\mu(C^c)+\mu(C\setminus K)<\frac{\delta}{2}+\frac{\delta}{2}=\delta,
\end{align*}
且 $f|_K$ 连续（因 $f|_C$ 连续且 $K\subseteq C$）。回到一般情形，对 $E$ 应用上述论证即得 $\mu(E\setminus K)<\delta$。
\end{proof}

% SUPPLEMENT: parent=theorem:ch02-Lusin; kind=motivation; 说明 Lusin 定理的意义、正则性被使用的两处位置、$\mu(E)<\infty$ 不可省的反例、连续性的相对含义以及与拟连续概念的关系
\begin{remark}\phantomsection\label{remark:ch02-Lusin}
\begin{enumerate}[(i)]
\item 引入动机。Lusin 定理给出"可测函数在去掉任意小测度集后连续"这一精确表述。它把可测性（集合层面的条件）与连续性（拓扑层面的条件）联系起来，是\cref{definition:ch02-拟连续}中拟连续概念的来源，也是\cref{proposition:ch02-Tietze逼近}中用连续函数逼近可测函数的依据。
\item 正则性的使用位置。第一步用内正则性把 $A$ 从内部用紧集 $K$ 逼近，第二步用正则性同时选取紧集 $K_n$ 与开集 $A_n$ 夹住 $f^{-1}(B_n)$，第四步再次用内正则性把闭集 $C$ 收紧为紧集 $K$。若 $\mu$ 不正则，这些逼近都不一定存在，结论一般失效。
\item $\mu(E)<\infty$ 不可省。取 $X=\mathbb{R}$、$\mu=\lambda$、$f=\chi_{\mathbb{Q}}$（Borel 可测）、$E=\mathbb{R}$。则 $\lambda(E)=\infty$。若存在紧集 $K\subseteq\mathbb{R}$ 使 $\lambda(\mathbb{R}\setminus K)<\delta$，则因 $K$ 有界，$\mathbb{R}\setminus K$ 含某个半直线，测度为 $\infty$，矛盾。故定理中"$\mu(E)<\infty$"不能去掉。证明中该假设用于 $\mu(A)\leqslant\delta/2<\infty$ 与 $\mu(C\setminus K)$ 的估计。
\item 连续性的含义。$f|_K$ 连续指相对于 $K$ 的相对拓扑连续，$K$ 可以没有内点（例如 Cantor 集）。定理只断言存在这样的紧集，不给出 $K$ 的构造；$K$ 依赖于 $E$ 与 $\delta$。
\item 与拟连续的关系。若 $\mu$ 有限且正则，则每个 $\mathscr B(X)$-可测函数都是 $\mu$-拟连续的（\cref{definition:ch02-拟连续}）：给定 $E\in\mathscr B(X)$ 与 $\delta>0$，对 $E$ 应用本定理得紧集 $K\subseteq E$，它特别为闭集，故满足拟连续性的要求。
\end{enumerate}
\end{remark}

\begin{definition}\label{definition:ch02-拟连续}
设 $X$ 为 Hausdorff 空间，$\mu$ 为 $\mathscr B(X)$ 上的测度。设 $E\in\mathscr B(X)$。称函数 $f:X\to\mathbb{R}$ 在 $E$ 中为 \textbf{$\mu$-拟连续}（或简称拟连续），若对任意 $\delta>0$，存在闭集 $C\subseteq E$ 使 $\mu(E\setminus C)<\delta$ 且限制 $f|_C$ 连续。
\end{definition}

% SUPPLEMENT: parent=definition:ch02-拟连续; kind=motivation; 说明拟连续概念的动机、$E\in\mathscr B(X)$ 与闭集 $C$ 的作用、相对连续性的含义，以及拟连续性不被几乎处处修改保持的边界
\begin{remark}\phantomsection\label{remark:ch02-拟连续}
\begin{enumerate}[(i)]
\item 引入动机。拟连续性是 Lusin 定理的抽象化：不要求函数本身连续，只要求它在去掉任意小测度集后连续。由\cref{theorem:ch02-Lusin}，有限正则测度下的 Borel 可测函数都是拟连续的；\cref{proposition:ch02-拟连续a.e.可测}则给出反向结论（至多差零测集）。
\item 条件的作用。$E\in\mathscr B(X)$ 使 $\mu(E\setminus C)$ 有意义；$C$ 取闭集（而非紧集）使定义不依赖空间的局部紧性；$\delta$ 任意小，故 $C$ 可以任意接近 $E$。连续性指 $f|_C$ 关于 $C$ 的相对拓扑连续。若把"闭集"换成"紧集"，则得到更强的概念，需要 $\mu$ 在 $E$ 上内正则才能保证二者一致。
\item 边界情形。拟连续性是相对于 $E$ 的性质：$f$ 在 $E$ 中拟连续不蕴含它在更大的集合中拟连续（$E$ 之外的取值不受控制）。拟连续性也不被 $\mu$-a.e. 相等保持：在零测集上修改 $f$ 可以破坏 $f|_C$ 的连续性，故拟连续性是函数本身（而非等价类）的性质。
\item 典型例子。$f=\chi_{(0,\infty)}$ 在 $\mathbb{R}$ 上关于 Lebesgue 测度拟连续：对 $\delta>0$，取 $C\triangleq\mathbb{R}\setminus(-\delta/2,\delta/2)$，则 $C$ 闭、$\lambda(\mathbb{R}\setminus C)=\delta$，且 $f|_C$ 为常值（在 $C\cap(0,\infty)$ 上为 $1$，在 $C\cap(-\infty,0]$ 上为 $0$），故连续。但 $f$ 在 $0$ 处不连续。
\end{enumerate}
\end{remark}

利用此定义，Lusin 定理可重述为：若 $\mu$ 有限且正则，则每个 $\mathscr B(X)$-可测函数均为 $\mu$-拟连续。下列结果表明逆命题（至多差零测集）成立。

\begin{proposition}\label{proposition:ch02-拟连续a.e.可测}
设 $X$ 为 Hausdorff 空间，$\mu$ 为 $\mathscr B(X)$ 上的有限正则测度。设 $f:X\to\mathbb{R}$ 为 $\mu$-拟连续。则存在 $\mathscr B(X)$-可测函数 $g:X\to\mathbb{R}$，使 $g=f$ 在 $X$ 上 $\mu$-a.e.。
\end{proposition}

% SUPPLEMENT: parent=proposition:ch02-拟连续a.e.可测; kind=proof; 原证明中"$f|_{C_n}=f\chi_{C_n}$ 连续"仅在 $C_n$ 的相对拓扑下成立，此处补出零延拓 $f\chi_{C_n}$ 的可测性论证、$\{C_n\}$ 非降化的验证与逐点收敛的证明
\begin{proof}
由\cref{definition:ch02-拟连续}，对每个 $n\in\mathbb{N}$，存在闭集 $C_n\subseteq X$ 使 $\mu(C_n^c)<1/n$ 且 $f|_{C_n}$ 连续。不妨设 $\{C_n\}$ 非增：否则以 $C_n'\triangleq\bigcup_{k=1}^nC_k$ 代替 $C_n$，它仍为闭集，$f|_{C_n'}$ 连续，且 $(C_n')^c\subseteq C_n^c$ 给出 $\mu((C_n')^c)<1/n$。令
\begin{align*}
C\triangleq\bigcup_{n=1}^\infty C_n,\qquad g\triangleq f\chi_C.
\end{align*}
则 $g=f$ 在 $C$ 上成立，且 $\mu(C^c)\leqslant\mu(C_n^c)<1/n$ 对一切 $n$ 成立，故 $\mu(C^c)=0$。

可测性。固定 $n\in\mathbb{N}$。因 $C_n$ 为闭集，其相对 Borel $\sigma$-代数等于 $\mathscr B(X)\cap C_n$；由 $f|_{C_n}$ 连续得：对任意 $B\in\mathscr B(\mathbb{R})$，
\begin{align*}
(f|_{C_n})^{-1}(B)\in\mathscr B(X)\cap C_n\subseteq\mathscr B(X).
\end{align*}
令 $g_n\triangleq f\chi_{C_n}$。则对任意 $B\in\mathscr B(\mathbb{R})$，
\begin{align*}
g_n^{-1}(B)=(f|_{C_n})^{-1}(B)\cup\bigl(C_n^c\ \text{若}\ 0\in B\bigr).
\end{align*}
其中 $(f|_{C_n})^{-1}(B)\in\mathscr B(X)$，而 $C_n^c\in\mathscr B(X)$。故 $g_n$ 为 $\mathscr B(X)$-可测。

逐点收敛。若 $x\in C$，则存在 $N$ 使 $x\in C_N$，从而 $n\geqslant N$ 时 $x\in C_n$，故 $g_n(x)=f(x)=g(x)$；若 $x\notin C$，则 $g_n(x)=0=g(x)$ 对一切 $n$ 成立。故 $g_n\to g$ 在 $X$ 上逐点。由\cref{corollary:ch02-可测运算}(ii)，$g$ 为 $\mathscr B(X)$-可测。因 $f=g$ 在 $C$ 上且 $\mu(C^c)=0$，结论成立。
\end{proof}

% SUPPLEMENT: parent=proposition:ch02-拟连续a.e.可测; kind=motivation; 说明与 Lusin 定理合起来给出的等价刻画、各假设在本方向中的实际使用情况、零延拓构造的作用以及拟连续性不被几乎处处修改保持的边界
\begin{remark}\phantomsection\label{remark:ch02-拟连续可测}
\begin{enumerate}[(i)]
\item 引入动机。本命题是 Lusin 定理的反向陈述：拟连续函数至多在一个零测集上与某个 Borel 可测函数不同。与\cref{theorem:ch02-Lusin}合起来，在 $\mu$ 有限且正则时给出
\begin{align*}
f\ \text{为}\ \mathscr B(X)\text{-可测}\quad\Longleftrightarrow\quad f\ \text{为}\ \mu\text{-拟连续（至多差一个零测集）}.
\end{align*}
\item 假设的实际使用。本方向的证明只用到 $C_n$ 为闭集（从而 $C_n^c\in\mathscr B(X)$）与 $\mu(C_n^c)<1/n$；$\mu$ 的有限性与正则性在此方向未被使用，它们是为\cref{theorem:ch02-Lusin}（另一方向）所需。因此本命题对任意 $\mathscr B(X)$ 上的测度成立。
\item 零延拓构造的作用。$f|_{C_n}$ 连续只保证它在 $C_n$ 上（相对拓扑）可测；要把它看成 $X$ 上的可测函数，须作零延拓 $g_n=f\chi_{C_n}$。零延拓在 $C_n$ 的边界上一般不连续，但可测性只要求原像可测，故 $g_n$ 可测。$\{C_n\}$ 非降化保证 $\{g_n\}$ 逐点收敛于 $g$。
\item 边界情形。结论中的 $g$ 在 $C^c$ 上的取值是任意的（$C^c$ 为零测集）；$g$ 依赖 $\{C_n\}$ 的选取，但任意两个这样的 $g$ 只在一个零测集上不同。拟连续性本身不被 $\mu$-a.e. 相等保持，故本命题不能加强为"$f$ 的每个等价类都有连续代表元"。
\end{enumerate}
\end{remark}

\begin{proposition}\label{proposition:ch02-Tietze逼近}\label{proposition:ch02-连续逼近}
设 $X$ 为局部紧 Hausdorff 空间，$\mu$ 为 $\mathscr B(X)$ 上的正则测度。设 $f:X\to\mathbb{R}$ 为 $\mu$-拟连续。则对任意开集 $A\subseteq X$（$\mu(A)<\infty$）与任意 $\delta>0$，存在紧集 $K\subseteq A$ 与函数 $\zeta\in C_c(X)$（$\operatorname{supp}\zeta\subseteq A$），使
\begin{align*}
\mu\bigl(\{f\neq\zeta\}\cap A\bigr)<\delta,\qquad
\|\zeta\|_\infty=\|f\|_K\triangleq\sup_{x\in K}|f(x)|.
\end{align*}
特别地，若 $f$ 在 $X\setminus A$ 上 $\mu$-a.e. 为零，则 $\mu(\{f\neq\zeta\})<\delta$。
\end{proposition}

% SUPPLEMENT: parent=proposition:ch02-Tietze逼近; kind=proof; 原证明引用 Lusin 定理、Tietze 延拓与 Urysohn 引理而未展开，此处补出由拟连续性选取紧集 $K$、Tietze 延拓的范数估计、截断函数的连续性与支撑估计以及测度估计
\begin{proof}
设 $A\subseteq X$ 为开集且 $\mu(A)<\infty$，并固定 $\delta>0$。

第一步：选取紧集。由 $f$ 的 $\mu$-拟连续性，存在闭集 $C\subseteq A$ 使 $\mu(A\setminus C)<\delta/2$ 且 $f|_C$ 连续。因 $C\in\mathscr B(X)$ 且 $\mu(C)\leqslant\mu(A)<\infty$，由 $\mu$ 的内正则性存在紧集 $K\subseteq C$ 使 $\mu(C\setminus K)<\delta/2$。于是
\begin{align*}
\mu(A\setminus K)\leqslant\mu(A\setminus C)+\mu(C\setminus K)<\delta,
\end{align*}
且 $f|_K$ 连续。

第二步：Tietze 延拓。$K$ 为紧集，故为 $X$ 的闭子集；$X$ 为局部紧 Hausdorff 空间，故正规。由 Tietze 延拓定理，存在 $h\in C(X)$ 使 $h|_K=f|_K$。

第三步：截断函数。由\cref{lemma:appA-Urysohn}，存在 $g\in C_c(X)$ 使 $g(X)\subseteq[0,1]$、$g|_K=1$、$\operatorname{supp}g\subseteq A$。令 $\eta\triangleq gh$。则 $\eta\in C_c(X)$，因为 $\operatorname{supp}\eta\subseteq\operatorname{supp}g$ 为紧集且含于 $A$；并且 $\eta|_K=h|_K=f|_K$。

第四步：范数截断。记 $M\triangleq\|f\|_K=\sup_{x\in K}|f(x)|<\infty$（$K$ 紧、$f|_K$ 连续）。定义
\begin{align*}
\zeta(x)\triangleq
\begin{cases}
\eta(x), & |\eta(x)|\leqslant M,\\
M\operatorname{sgn}\eta(x), & \text{否则}.
\end{cases}
\end{align*}
则 $\zeta=\varphi\circ\eta$，其中 $\varphi(t)\triangleq t$（$|t|\leqslant M$）、$\varphi(t)\triangleq M\operatorname{sgn}t$（$|t|>M$）为连续函数；故 $\zeta\in C(X)$，且 $\operatorname{supp}\zeta\subseteq\operatorname{supp}\eta\subseteq A$ 紧，即 $\zeta\in C_c(X)$。由 $|\zeta|\leqslant M$ 得 $\|\zeta\|_\infty\leqslant M$。

第五步：范数等式与测度估计。对 $x\in K$ 有 $\eta(x)=f(x)$ 与 $|f(x)|\leqslant M$，故 $\zeta(x)=\eta(x)=f(x)$。因 $K$ 紧、$|f|$ 在 $K$ 上连续，存在 $x_0\in K$ 使 $|f(x_0)|=M$，故 $|\zeta(x_0)|=M$，从而 $\|\zeta\|_\infty=M=\|f\|_K$。此外，$x\in K$ 时 $f(x)=\zeta(x)$，故
\begin{align*}
\{f\neq\zeta\}\cap A\subseteq A\setminus K,
\end{align*}
于是 $\mu(\{f\neq\zeta\}\cap A)\leqslant\mu(A\setminus K)<\delta$。最后，若 $f$ 在 $X\setminus A$ 上 $\mu$-a.e. 为零，则由 $\operatorname{supp}\zeta\subseteq A$ 得 $\zeta=0$ 在 $X\setminus A$ 上成立，故
\begin{align*}
\{f\neq\zeta\}\subseteq\bigl((X\setminus A)\cap\{f\neq0\}\bigr)\cup\bigl(A\cap\{f\neq\zeta\}\bigr),
\end{align*}
右端第一个集合为 $\mu$-零测集，第二个集合的测度小于 $\delta$，故 $\mu(\{f\neq\zeta\})<\delta$。
\end{proof}

% SUPPLEMENT: parent=proposition:ch02-Tietze逼近; kind=motivation; 说明用紧支连续函数逼近的用途、局部紧与 $A$ 的有限测度的作用、范数等式的意义，以及 $\mu(\{f\neq\zeta\})$ 必须限制在 $A$ 上的反例
% SUPPLEMENT: parent=proposition:ch02-连续逼近; kind=motivation; 别名标签，与 proposition:ch02-Tietze逼近 共用下一条注记
\begin{remark}\phantomsection\label{remark:ch02-Tietze逼近}
\begin{enumerate}[(i)]
\item 引入动机。本命题给出定量的 Lusin 型逼近：用紧支连续函数 $\zeta$ 在测度意义下逼近拟连续函数 $f$，并保持 $\sup$ 范数。它是 $C_c(X)$ 在 $L^p(X)$ 中稠密（\cref{lemma:ch02-Cc稠密}）以及 Young 测度理论中构造测试函数的标准工具。
\item 假设的作用。$X$ 局部紧 Hausdorff 保证 Urysohn 引理与 Tietze 延拓可用，并使"紧支连续函数"有足够多；$\mu(A)<\infty$ 保证 $K\subseteq A$ 可从内部用紧集逼近且 $\mu(A\setminus K)$ 可估计，同时 $M=\|f\|_K<\infty$；拟连续性给出第一步的闭集 $C$；$\mu$ 的正则性用于把 $C$ 收紧为紧集 $K$。
\item 范数等式的作用。$\|\zeta\|_\infty=\|f\|_K$ 使逼近不放大函数的振幅，这在后续把 $\zeta$ 代入线性泛函或积分时控制估计。若不作第四步的截断，$\eta=gh$ 的 $\sup$ 范数可能大于 $M$。
\item $\mu(\{f\neq\zeta\})$ 必须限制在 $A$ 上。$\zeta$ 在 $A$ 外恒等于 $0$（因 $\operatorname{supp}\zeta\subseteq A$），故若 $f$ 在 $A$ 外不几乎处处为零，则 $\{f\neq\zeta\}$ 含 $A$ 外测度为正的部分。反例：取 $X=\mathbb{R}$、$\mu=\lambda$、$f\equiv1$、$A=(0,1)$、$\delta=1/2$，则任何满足 $\operatorname{supp}\zeta\subseteq A$ 的 $\zeta$ 在 $\mathbb{R}\setminus A$ 上为零，故 $\lambda(\{f\neq\zeta\})=\infty$。本命题的标准用法是 $f$ 在 $A$ 外 $\mu$-a.e. 为零（例如 $f=\chi_E$，$E\subseteq A$），此时补充断言给出 $\mu(\{f\neq\zeta\})<\delta$。
\end{enumerate}
\end{remark}

\begin{remark}\phantomsection\label{remark:ch02-Lusin对可测函数}
由 Lusin 定理，\cref{proposition:ch02-Tietze逼近}对 $\mathscr B(X)$-可测函数也成立。事实上，设 $f:X\to\mathbb{R}$ 为 $\mathscr B(X)$-可测，$A\subseteq X$ 为开集且 $\mu(A)<\infty$，$\delta>0$。由\cref{theorem:ch02-Lusin}（取 $E=A$）存在紧集 $K\subseteq A$ 使 $\mu(A\setminus K)<\delta/2$ 且 $f|_K$ 连续；再重复\cref{proposition:ch02-Tietze逼近}证明的第二步至第五步，即得满足所述性质的 $\zeta\in C_c(X)$。
\end{remark}

对容量有类似概念与结果。

\begin{definition}\label{definition:ch02-C拟连续}
设 $(X,d)$ 为度量空间，$\mathscr F\subseteq\mathscr P(X)$ 如\cref{definition:ch01-容量}所述且对补集封闭，$C:\mathscr F\to[0,\infty]$ 为容量。称函数 $f:X\to\mathbb{R}$ 在 $X$ 上为 \textbf{$C$-拟连续}，若对任意 $\delta>0$，存在子集 $E\subseteq X$（$E^c\in\mathscr F$）使 $C(E^c)<\delta$ 且限制 $f|_E$ 连续。
\end{definition}

% SUPPLEMENT: parent=definition:ch02-C拟连续; kind=motivation; 说明容量拟连续的动机、$E^c\in\mathscr F$ 与相对连续性的作用、与测度拟连续的差别，以及外容量情形的边界
\begin{remark}\phantomsection\label{remark:ch02-外容量开集}
\begin{enumerate}[(i)]
\item 引入动机。$C$-拟连续是拟连续概念的容量版本：它不要求 $E$ 闭或紧，只要求其补集的容量任意小。它用于研究 Sobolev 函数、容量零集与精细性质，例如 $2$-拟连续代表元的存在性（见\cref{definition:ch02-拟连续}）。
\item 条件的作用。$E^c\in\mathscr F$ 使 $C(E^c)$ 有定义；连续性指 $f|_E$ 关于 $E$ 的相对拓扑连续。与测度拟连续不同，这里允许 $E$ 为任意子集，从而不必验证 $E$ 可测；容量只作用于 $E^c$。
\item 外容量情形可设 $E^c$ 为开集。设 $C$ 为外容量。给定 $\delta>0$，取 $E\subseteq X$ 使 $C(E^c)<\delta$ 且 $f|_E$ 连续。由外容量定义 $C(E^c)=\inf\{C(A)\mid A\supseteq E^c,\ A\ \text{开}\}$，故存在开集 $A\supseteq E^c$ 使 $C(A)<\delta$。令 $E'\triangleq A^c$。则 $E'$ 为闭集、$E'\subseteq E$，故 $f|_{E'}$ 为 $f|_E$ 的限制，连续；又 $(E')^c=A$ 为开集且 $C((E')^c)<\delta$。故在定义中不妨设 $E^c$ 为开集，即 $E$ 为闭集。
\item 边界情形。若 $\mathscr F$ 不对补集封闭，则 $C(E^c)$ 无意义，定义需另行解释。$C$-拟连续性是函数本身的性质，不被 $C$-q.e. 相等保持（与测度情形相同）。
\end{enumerate}
\end{remark}

\begin{proposition}\label{proposition:ch02-容量拟连续闭}
设 $(X,d)$ 为度量空间，$\mathscr F\subseteq\mathscr P(X)$ 如\cref{definition:ch01-容量}所述且对补集封闭，$C:\mathscr F\to[0,\infty]$ 为容量。设 $f,f_n:X\to\mathbb{R}$（$n\in\mathbb{N}$）。设每个 $f_n$ 为 $C$-拟连续，且 $f_n\to f$ 依容量。则 $f$ 为 $C$-拟连续。
\end{proposition}

% SUPPLEMENT: parent=proposition:ch02-容量拟连续闭; kind=proof; 原证明中 $C(G^c)$ 的估计笔误与"一致极限保持连续"的使用均已补全，此处给出 $E$、$F$、$G$ 的容量估计与连续性论证的完整版本
\begin{proof}
固定 $\delta>0$。由 $C$-拟连续性，对每个 $n\in\mathbb{N}$ 存在子集 $E_n\subseteq X$ 使 $C(E_n^c)<\delta/2^{n+1}$ 且 $f_n|_{E_n}$ 连续。令
\begin{align*}
E\triangleq\bigcap_{n=1}^\infty E_n.
\end{align*}
则 $E^c=\bigcup_{n=1}^\infty E_n^c\in\mathscr F$，且由可数次可加性
\begin{align*}
C(E^c)\leqslant\sum_{n=1}^\infty C(E_n^c)\leqslant\frac{\delta}{2}.
\end{align*}
又 $E\subseteq E_n$，故 $f_n|_E$ 为 $f_n|_{E_n}$ 的限制，连续（$n\in\mathbb{N}$）。

由\cref{proposition:ch02-容量收敛关系}(ii)，存在子列 $\{f_{n_k}\}\subseteq\{f_n\}$ 与 $F\subseteq X$ 使 $C(F^c)<\delta/2$ 且 $f_{n_k}\to f$ 在 $F$ 上一致。令
\begin{align*}
G\triangleq E\cap F.
\end{align*}
则 $G^c=E^c\cup F^c\in\mathscr F$，且
\begin{align*}
C(G^c)\leqslant C(E^c)+C(F^c)<\delta.
\end{align*}
对每个 $k$，$f_{n_k}|_G$ 为 $f_{n_k}|_{E}$ 的限制，连续；又 $f_{n_k}\to f$ 在 $G$ 上一致。连续函数列的一致极限连续，故 $f|_G$ 连续。由 $\delta$ 的任意性，$f$ 为 $C$-拟连续。
\end{proof}

% SUPPLEMENT: parent=proposition:ch02-容量拟连续闭; kind=motivation; 说明容量拟连续在依容量收敛下的稳定性、可数次可加性与子列拟一致收敛的作用，以及仅假设 $C$-q.e. 收敛时结论失效的边界
\begin{remark}\phantomsection\label{remark:ch02-容量拟连续闭}
\begin{enumerate}[(i)]
\item 引入动机。本命题说明 $C$-拟连续函数类在依容量收敛下封闭。它是容量框架下与\cref{proposition:ch02-收敛关系}(ii) 平行的结果，用于证明 Sobolev 空间中拟连续代表元的唯一性与稳定性。
\item 假设的作用。每个 $f_n$ 的 $C$-拟连续性提供 $E_n$；依容量收敛通过\cref{proposition:ch02-容量收敛关系}(ii) 提供 $C$-拟一致收敛的子列与集合 $F$；$C$ 的可数次可加性用于估计 $C(E^c)$ 与 $C(G^c)$。三者缺一不可。
\item 边界情形。若把"依容量收敛"减弱为"$C$-q.e. 收敛"，结论不再成立：$C$-q.e. 收敛不提供一致的容量估计，无法把 $E^c$ 与 $F^c$ 的容量同时控制。此外，$E$ 与 $F$ 只要求为子集，不要求闭或紧，这与测度拟连续情形（要求闭集）不同。
\end{enumerate}
\end{remark}

\end{document}
